\documentclass[11pt,a4paper,oneside]{amsbook}
\usepackage[T1]{fontenc}
\usepackage[utf8]{inputenc}
\usepackage{lmodern}
\usepackage[american]{babel}
\usepackage{microtype}
\usepackage{geometry}
\geometry{margin=1.1in}
\makeatletter
\def\input@path{{../}}
\makeatother
\usepackage{sga1-en}
\usepackage[hidelinks]{hyperref}

\addto\captionsamerican{\renewcommand{\chaptername}{Expos\'e}}

\title[SGA~1, Expos\'e V (English draft)]{The fundamental group: generalities\\
{\normalsize SGA~1, Expos\'e V\\English draft}}
\author{A.\ Grothendieck}
\date{}
\hypersetup{
  pdftitle={SGA 1, Expose V: The fundamental group: generalities (English draft)},
  pdfauthor={A. Grothendieck},
  pdfsubject={Unofficial English translation of SGA 1, Expose V}
}

\begin{document}
\frontmatter
\maketitle

\chapter*{About this draft}
This is an unofficial English translation of Expos\'e~V,
``Le groupe fondamental: g\'en\'eralit\'es'', from SGA~1. The source volume is
A.\ Grothendieck and M.\ Raynaud,
\textit{Rev\^etements \'etales et groupe fondamental} (SGA~1),
S\'eminaire de g\'eom\'etrie alg\'ebrique du Bois Marie, 1960--61.
It follows the slightly corrected SMF recomposition,
\href{https://arxiv.org/abs/math/0206203v2}{arXiv:math/0206203v2}.

This draft contains the whole expos\'e, including proofs, footnotes,
diagrams, the updating remarks added in 2003 by M.~Raynaud (between bold
brackets and marked (MR)) where the expos\'e has them, and its bibliography.
Statement numbering and mathematical notation follow the source.
References to other expos\'es print the source's numbers. Apparent
misprints in the source have been retained; they are recorded in the
accompanying README, separately from the translated text. Scholarly
proofreading remains outstanding.

The translator's contribution is licensed under
the MIT License.
The French original remains copyright of its authors and publishers.

\tableofcontents
\mainmatter
\setcounter{chapter}{4}
\chapter{The fundamental group: generalities}
\label{V}
% original p. 105

\setcounter{section}{-1}
\section{Introduction}

The present Seminar is the continuation of the 1960 Seminar. We
refer to the latter by symbols such as I~\Ref{I.9.7},
which means: S\'eminaire de G\'eom\'etrie Alg\'ebrique,
Expos\'e~I, \No9.7. The numbers of the 1961 expos\'es will follow
those of 1960. We refer to the \'El\'ements de
G\'eom\'etrie Alg\'ebrique of Dieudonn\'e-Grothendieck by
symbols such as (EGA~I 8.7.3).

The present expos\'e summarizes (with slight
complements) the last expos\'es of 1960, which had not
been written up.

As in 1961, we shall as a general rule restrict ourselves
to locally noetherian preschemes, although this restriction is
often inessential. In
expos\'e \Ref{VI} we shall take for granted the theory of faithfully flat descent,
summarized in S\'eminaire Bourbaki \No190. If appropriate,
we shall give a more detailed exposition of it in a
later expos\'e\footnote{\Cf Exp\ptbl \Ref{VI} and Exp\ptbl \Ref{VIII}}, once
the reader has had the opportunity to convince himself of the usefulness of
this technique for the theory of the fundamental group.

\section{Prescheme with a finite group of operators,
quotient prescheme}
\label{V.1}

Let $X$ be a prescheme, $G$ a finite group operating on~$X$ by automorphisms, on the right to fix ideas. If $X$
is affine with ring $A$, $G$ therefore operates by automorphisms on
the left on~$A$.

For every prescheme $Z$, $G$ operates on the left on
the set $\Hom(X,Z)$,
one can therefore consider the set
$$
\Hom(X,Z)^G
$$
of morphisms invariant under $G$. It depends functorially on
$Z$,
one can ask whether this functor is ``representable'',
\ie isomorphic to a functor $Z\mto\Hom(Y,Z)$.
% original p. 106
This means that one can find a prescheme $Y$, and a
morphism invariant under $G$
$$
p\colon X\to Y
$$
such that for every $Z$, the corresponding map $g\mto gp$
$$
\Hom(Y,Z)\to \Hom(X,Z)^G
$$
is bijective. One then says that $(Y,p)$ is a
\emph{quotient prescheme}
of~$X$
by~$G$
(it is determined up to unique isomorphism).

\begin{proposition}
\label{V.1.1}
Let $A$ be a ring on which the finite group $G$ operates on
the left, $B=A^G$ the subring of invariants of~$A$, $X=\Spec(A)$ and
$Y=\Spec(B)$, $p\colon X\to Y$ the canonical morphism (obviously
invariant under $G$). Then
\begin{enumerate}
\item[(i)] $A$ is integral over~$B$, \ie $p$ is an
\emph{integral} morphism.
\item[(ii)] The morphism $p$ is surjective, its fibers are the
orbits of~$G$, the topology of~$Y$ is the quotient of that of~$X$.
\item[(iii)] Let $x\in X$, $y=p(x)$, $G_x$ the stabilizer of~$x$; then
$\kres(x)$ is a quasi-Galois algebraic extension of
$\kres(y)$ and the canonical map of~$G_x$ into the group
$\Gal(\kres(x)/\kres(y))$ of $\kres(y)$-automorphisms of
$\kres(x)$
is
surjective.
\item[(iv)] $(Y,p)$ is a quotient prescheme of~$X$ by $G$.
\end{enumerate}
\end{proposition}

Statements
(i), (ii), (iii)
are well known in commutative
algebra\footnote{\Cf N. Bourbaki, Alg. Comm. Chap\ptbl 5, \S 1 and \S
2, th\ptbl 2.} and are included only as a reminder, except for the assertion
about the topology, which comes from the following general fact,
an easy consequence of the Cohen-Seidenberg theorem: an
integral morphism is closed (\ie transforms closed sets into
closed sets). Let us note at once:

\begin{corollary}
\label{V.1.2}
Under the preceding conditions, the natural homomorphism
$$\cal{O}_Y{\to} p_*(\cal{O}_X)^G$$ is an isomorphism.
\end{corollary}

This follows at once from the formula
$$
(S^{-1}A)^G = S^{-1}(A^G)
$$
valid for every multiplicatively stable subset $S$ of~$B=A^G$
(a formula which extends to modules, and which can be stated more generally for
a change of base $A\to A'$ which is \emph{flat}),
applied
to the case where $S$ is generated by an element $f$ of~$B$.

Assertion (ii) and cor\ptbl \Ref{V.1.2} easily imply (iv); more
generally, one will have the following:

\begin{proposition}
\label{V.1.3}
Let
% original p. 107
$X$ be a prescheme with a finite group of automorphisms
$G$, $p\colon{}X\to{}Y$ an invariant affine morphism such that
$\cal{O}_Y\isomto p_*(\cal{O}_X)^G$. Then the conclusions
\textup{(i), (ii), (iii), (iv)}
of~\Ref{V.1.1} are still valid.
\end{proposition}

Indeed, for
(i), (ii), (iii),
one may assume $Y$, hence $X$, affine, and
if $B,A$ are their rings, the hypothesis implies $B=A^G$, it
suffices to apply~\Ref{V.1.1}. For (iv), one uses (ii) and
$\cal{O}_Y=p_*(\cal{O}_X)^G$.

\begin{corollary}
\label{V.1.4}
Under the conditions of~\Ref{V.1.3}, for every open set $U$ of~$Y$, $U$
is a quotient of~$X|U=p^{-1}(U)$ by $G$.
\end{corollary}

Indeed, $p^{-1}(U)\to{}U$ induced by $p$ satisfies the same
hypotheses as~$p$.

If now $X$ is a $Z$-prescheme and the operations
of~$G$ are $Z$-automorphisms, then by (iv) $Y$ is a
$Z$-prescheme. That said:

\begin{corollary}
\label{V.1.5}
For $X$ to be affine \resp separated over~$Z$, it is necessary and
sufficient that~$Y$ be so. If $X$ is of finite type over~$Z$, it is
\emph{finite} over~$Y$; if moreover $Z$ is locally noetherian, $Y$
is of finite type over~$Z$.
\end{corollary}

Since $X$ is affine and a fortiori separated over~$Y$, if $Y$ is
affine \resp separated over~$Z$, so
is~$X$. Conversely, suppose $X$ affine over~$Z$, let us prove that
$Y$ is: thanks to~\Ref{V.1.4} one may assume $Z$ affine, and
one is reduced to proving that if $X$ is affine, $Y$~is, which
follows from the explicit determination of~$Y$ as
$\Spec(A)^G$ made in~\Ref{V.1.1}. Likewise, since
$p\colon{}X\to{}Y$ is integral, hence universally closed, and
surjective, it follows that if $X$ is separated over~$Z$, so is
$Y$ (a lemma to be singled out!); indeed, in the diagram
$$
\xymatrix@C=1.5cm{ X \times_ZX\ar[r]^-{p\times_Zp} & Y\times_ZY \\ X
\ar[u]^{\Delta_{X/Z}}\ar[r]^{p} & Y \ar[u]_{\Delta_{Y/Z}} }
$$
the morphism $X\times_ZX\to Y\times_ZY$ is closed, hence transforms
the (closed) diagonal of~$X\times_ZX$ into a closed subset of
$Y\times_ZY$, which moreover is none other than the diagonal of the
latter since $p$ is surjective.
--- If $X$ is of finite type over~$Z$, it is a fortiori so over~$Y$, hence
it is finite over~$Y$ (since it is already integral over~$Y$). Suppose moreover $Z$ locally noetherian, let us prove that
$Y$ is of finite type over~$Z$. Thanks to~\Ref{V.1.4} one may
assume $Z$ affine. Since the topological space $X$ is quasi-compact
and $p\colon X\to Y$ is surjective, $Y$ is
likewise
quasi-compact, hence a finite union of
% original p. 108
affine open sets, and by~\Ref{V.1.4} one is reduced to the case where $Y$ is
affine, hence $X$ affine. But then the ring $A$ of~$X$ is an
algebra of finite type over the ring $C$ of~$Z$, which is
noetherian, and it is known that $B=A^G$ is then likewise
an algebra of finite type over~$C$ (for $A$ will be integral, hence
finite over a subalgebra $B'$ of~$B$ of finite type over~$C$, hence,
since $B'$ is noetherian, $B$ is likewise
finite
over~$B'$, hence of finite type over~$C$).

\begin{corollary}
\label{V.1.6}
For $X$ to be affine \resp a scheme, it is necessary and sufficient that $Y$ be
so.
\end{corollary}

\begin{definition}
\label{V.1.7}
Let
% original p. 109
$X$ be a prescheme on which a finite group $G$ operates
on the right. One says that $G$ \emph{operates admissibly}
if there exists a morphism $p\colon X\to Y$ having the properties
of~\Ref{V.1.3} (which implies that $X/G$ exists and is isomorphic
to~$Y$).
\end{definition}

\begin{proposition}
\label{V.1.8}
Let $X$ be a prescheme on which the finite group $G$ operates
on the right. For $G$ to operate admissibly, it is necessary
and sufficient that $X$ be a union of affine open sets invariant under
$G$, or again that every orbit of~$G$ in $X$ be contained
in an affine open set.
\end{proposition}

This last condition is obviously implied by the
first, and in its turn it implies it; for let $T$ be an
orbit of~$G$, $U$ an affine open set containing it; the intersection
of the transforms of~$U$ by the
$g$ in $G$
is then an open set $U'$ stable under~$G$, containing $T$ and contained in
the affine open set $U$. Since in $U$ every finite subset has a
fundamental system of affine open neighborhoods, there exists an
affine open neighborhood $V$ of~$T$ contained in $U'$. These
transforms by the
$g$ in $G$
are therefore affine and contained in $U'$, which is \emph{separated},
hence their intersection $U''$ is an affine open set which is invariant
under $G$ and contains $T$.
--- This being established, the condition considered in~\Ref{V.1.8} is
\emph{necessary}, for one will take the inverse images $X_i$
of affine open sets $Y_i$ covering~$Y$. It is sufficient, for one
can then by~\Ref{V.1.1} construct the quotients $Y_i=X_i/G$; in
each $Y_i$ the image of~$X_i\cap X_j$ is an open set $Y_{ij}$
which is identified with $X_{ij}/G$ by~\Ref{V.1.4}, in particular one
deduces from this isomorphisms $Y_{ij}\isomto Y_{ji}$ making it possible to
glue the $Y_i$ together to construct $Y$. Serre prefers to
construct directly the quotient topological space $Y$ of~$X$ by
$G$, to put on it the sheaf $p_*(\cal{O}_X)^G$ and to verify that
$Y$ becomes a prescheme and that one is then under the
conditions of~\Ref{V.1.3}.

\setcounter{subsection}{6}

\begin{corollary}
\label{cor:V.1.7}
If $G$ operating on~$X$ is admissible, the same holds for
every subgroup $H$ of~$G$ (hence $X/H$ exists).
\end{corollary}

This can also be verified directly on the
situation~\Ref{V.1.3}, by noting that one can always assume $X$
affine over some $Z$ and the $s\in G$ operating by $Z$-automorphisms
(one takes for example $Z=Y$); indeed one has:

\begin{corollary}
\label{cor:V.1.8}
Suppose $X$ affine over~$Z$, and the operations of~$G$ to be
$Z$-auto\-mor\-phisms. Then $G$ operates on~$X$
admissibly. If $X$ is defined by a quasi-coherent sheaf
$\cal{A}$ of algebras, $Y$ is defined by the sheaf
$\cal{A}^G$ of invariants of~$\cal{A}$ under~$G$.
\end{corollary}

\begin{proposition}
\label{V.1.9}
Suppose that $G$ operates admissibly on~$X$, and that
$X/G=Y$ is a prescheme over~$Z$. Consider a
change-of-base morphism $Z'\to Z$, set $X'=X\times_Z Z',
Y'=Y\times_Z Z'$, so that $G$ still operates by transport of
structure on~$X'$, the morphism $p'\colon X'\to Y'$ being
invariant. If $Z'$ is \emph{flat} over~$Z$, then $p'$ still satisfies
the hypotheses
of~\Ref{V.1.3},
\ie $\cal{O}_Y'\to
p_*'(\cal{O}_{X'})^G$ is an isomorphism ($p'$ being affine in any
case). Hence $G$ operates admissibly on~$X'$, and $(X/G)\times_Z Z'\thickapprox (X\times_Z Z')/G$.
\end{proposition}

One may obviously assume $Z=Y$, and one is reduced to the case where
moreover $Y$ and $Y'$ are affine. One must show that if $B$ is the
subring of invariants of~$G$ operating in $A$, and if $B'$ is
an algebra over~$B$ which is flat over~$B$, then $B'$ is the
subalgebra of invariants of~$A'=A\otimes_B B'$. This is
immediate, for the exact sequence
$$
0 \to B \lto{i} A \lto{j} A^{(G)}
$$
(where the last term means a power of~$A$, and where
$j(x)$ is the system of the $s\cdot x-x$, $s\in G$) remains exact under
tensoring with
$B'$.

One should take care that the flatness hypothesis was
essential for the validity of the result; in particular, if
$Y'$ is a closed subprescheme of
$Y$
(for example even a closed point of
$Y$),
$X'$ its inverse image in $X$,
then $Y'$ \emph{is not identified} in general with
$X'/G$. We shall see that it is nevertheless so if $X$ is
\'etale over~$Y$.

To finish, let us give a formalism as convenient as it is trivial. Let $Y$
be a prescheme. Since in the category of
preschemes direct sums exist, one can for every
set $E$ consider the prescheme sum of a family
$(Y_i)_{i\in E}$ of preschemes all identical to $Y$; this
prescheme will be denoted $Y\times E$.
% original p. 110
It is characterized by the formula
\begin{equation*}
\label{eq:V.1.*}
\tag{$*$} \Hom(Y\times E,Z) = \Hom(E,\Hom(Y,Z))
\end{equation*}
where the second $\Hom$ obviously denotes the set
of maps from the set $E$ into the set $\Hom(Y,Z)$. One has
a canonical morphism
$$
Y\times E\to Y
$$
making~$Y\times E$ a prescheme over~$Y$. Since
fibered products commute with direct sums (in the
category of preschemes) one will have, if $Y$ is a
prescheme over some other $Z$, for a change of base $Z'\to
Z$:
$$
(Y\times E)\times_Z Z'=(Y\times_Z Z')\times E
$$
(a formula especially useful if $Z=Y$). On the other hand, one concludes
trivially from the definition
$$
(Y\times E)\times F=Y\times(E\times F)=(Y\times E)\times_Y(Y\times F)
$$
(the last formula however resulting from the
commutativity pointed out above).

For fixed $Y$, one can regard $Y\times E$ as a functor in
$E$, with values in preschemes over~$Y$, a functor which
commutes with finite products by the preceding formula
(which makes it possible for example to associate with every ordinary group $G$ a
group scheme $Y\times G$ over~$Y$, which will be
finite over~$Y$ if $Y$ is, etc.). More generally, this
functor is ``left exact'', but we shall not need
to use this here. This functor also commutes trivially with direct
sums, and it is also ``right exact'', as one sees
at once from the defining formula \eqref{eq:V.1.*}. In
particular, if the finite group $G$ operates on the right in
the set $E$, then it operates on the right in $Y\times E$, and
one has
$$
(Y\times E)/G = Y\times (E/G)
$$
where in fact the quotient of the left-hand side satisfies the conditions
of~\Ref{V.1.3} (this is immediate).

\section{Decomposition and inertia groups. \'Etale case}
\label{V.2}
Let $G$ be a
finite group acting on the right on the prescheme $X$. If
$x\in X$, one calls
% original p. 111
\emph{decomposition group of~$x$}
the stabilizer $G_d(x)$ of~$x$. This group acts canonically
(on the left) on the residue field
$\kres(x)$, and the set
of elements of~$G_d(x)$ that act trivially is
called the \emph{inertia group}
of~$x$, denoted~$G_i(x)$.

Suppose that $G$ acts on~$X$ in an admissible way and that
$Y$ is a prescheme over a prescheme
$Z$. Let us fix a $z\in Z$, and an algebraically
closed extension $\Omega$
of~$\kres(z)$ having a transcendence degree
greater than that of the $\kres(x)/\kres(z)$, where $x$ is a
point of~$X$ over~$z$. One can regard $\Spec(\Omega)$ as
a $Z$-scheme, and the points of~$X$ with values in $\Omega$
correspond to the homomorphisms
of~$\kres(z)$-algebras $\kres(x)\to\Omega$,
where $x$ is a point of~$X$ over~$z$;
since $\Omega$ has been taken large enough, every point $x$ of~$X$
over~$z$ is the locality of a point of~$X$ with values
in $\Omega$. If $X(\Omega)$ and $Y(\Omega)$ denote
respectively the set of points of~$X$ and $Y$ with values in
$\Omega$, one has a natural map
$$
X(\Omega)\to Y(\Omega),
$$
on the other hand, $G$ acts on~$X(\Omega)$ and the
preceding map is invariant under $G$. This being so, the
conclusions (ii) and (iii) of~\Ref{V.1.3} can also be interpreted
as follows: \emph{the preceding map is surjective and
identifies $Y(\Omega)$ with the quotient $X(\Omega)/G$. Moreover, if $x$ is
the locality of~$a\in X(\Omega)$, then the stabilizer of~$a$
in $G$ is none other than the inertia group $G_i(x)$}. All this is
moreover true without assuming $\Omega$ ``large enough''; this
last hypothesis serves only to ensure that one can
characterize the inertia group of every element of~$X$
over~$z$ as a ``geometric'' stabilizer. One
concludes from this, for instance, immediately:
\begin{proposition}
\label{V.2.1}
Let us make an extension of the base $Z'\to Z$, whence
$X'=X\times_ZZ'$. Let $x'$ be a point of~$X'$, $x$ its image in $X$;
then one has $G_i(x)=G_i(x')$.
\end{proposition}

It suffices, in the considerations above, to take for
$\Omega$ a sufficiently large extension
of~$\kres(z')$
(where $z,z'$ are the images of~$x,x'$ in $Z,Z'$).

\begin{proposition}
\label{V.2.2}
Under the conditions of~\Ref{V.1.3}, suppose $Y$ locally
noetherian, $X$ finite over~$Y$. Let $H$ be a subgroup of~$G$,
consider $X'=X/H$ (\cf \Ref{cor:V.1.7}), let
% original p. 112
$x\in X$, $x'$ its image in $X'$ and $y$ its image in $Y$.
\begin{enumerate}
\item[(i)] If $H\supset G_d(x)$, then the homomorphism
$\cal{O}_y\to\cal{O}_{x'}$ induces an isomorphism on the
completions.
\item[(ii)] If $H\supset G_i(x)$, then the homomorphism $\cal{O}_y\to
\cal{O}_{x'}$ is \'etale, \ie $X'$ is \'etale over~$Y$ at $x'$.
\end{enumerate}
\end{proposition}
Let $Y_1=\Spec(\widehat{\cal{O}_y})$; let us make the change of base
$Y_1\to Y$; one finds an $X_1=X\times_Y Y_1$ finite over~$Y_1$, on
which $G$ acts, the quotient being $Y_1$
by~\Ref{V.1.9}. Let~$y_1$ be the unique point of~$Y_1$ over~$y$;
since
$\kres(y)=\kres(y_1)$,
it follows that the fiber of~$X$ at~$y$
is isomorphic to that of~$X_1$ at $y_1$, whence a unique point
$x_1$ of~$X_1$ over~$x$. Moreover, by~\Ref{V.1.9} one will have
$X_1/H=X'_1=X'\times_Y Y_1$; let $x'_1$ be the image of~$x_1$ in $X'_1$;
it lies over~$x'$, and one easily checks ($X'$ being
of finite type over~$Y$) that the homomorphism
$\cal{O}_{x'}\to\cal{O}_{x'_1}$ induces an isomorphism on the
completions. Hence one is reduced to the case where $Y$ is the
spectrum of a complete local ring, say $B$, hence~$X$ the spectrum of a
ring $A$ finite over~$B$, composed of a finite number of local rings~$A_x$ corresponding
to the points $x_i$ of~$X$ over~$Y$. If $A_0$ corresponds to $x=x_0$,
then $A$ is identified with the ring $\Hom_{G_d}(G,A_0)$ of the functions
$f\colon G\to A_0$ such that $f(st)=sf(t)$ for $s\in G_d$, the
operations of~$G$ on these functions being defined by
$(uf)(t)=f(tu)$. One sees therefore that if $H$ is an arbitrary subgroup
of~$G$, then $A^H$ is the ring of the functions $f\colon G\to A_0$
such that
$$
f(stu)=sf(t)\qquad\text{for}\ s\in G_d,\ u\in H
$$
hence it is a semi-local ring whose local components
correspond to the double cosets $G_daH$ in $G$, the double
coset defined by $a\in G$ corresponding (thanks to
the map $f\mto f(a)$) to the subring $A_0^{H(a)}$ of~$A_0$,
where $H(a)=G_d\cap aHa^{-1}$. Moreover, the local component of
$A^H$ corresponding to the image $x'$ of~$x$ is also the one
corresponding to the double coset $G_dH$ of the neutral
element; its local component is therefore $A_0^{G_d\cap H}$. If therefore
$G_d\subset H$,
one finds $A_0^{G_d}=A^G=B$, which proves (i). To
prove (ii), one can, by passing to a suitable finite flat
extension of~$A$, and using~\Ref{V.2.1}, reduce to the case where
the residue extension
$\kres(x)/\kres(y)$
is trivial. But
then $G_i(x)=G_d(x)$, and one is reduced to the preceding case.

\begin{corollary}
\label{V.2.3}
Under
% original p. 113
the conditions of~\Ref{V.2.2}, suppose $G_i(x)=(e)$; then X
is \'etale over~$Y$ at $x$. Hence if $G_i(x)=(e)$ for every $x \in
X$, then $X \to Y$ is an \'etale morphism.
\end{corollary}

There is a partial converse:

\begin{corollary}
\label{V.2.4} Suppose $X$ connected and the group $G$ faithful
on~$X$. For $p \colon X \to Y=X/G$ to be \'etale, it is necessary and
sufficient that the inertia groups of the points of~$X$ be reduced
to the neutral element. If this is so, $G$ is identified with the
group of all $Y$-automorphisms of the $Y$-scheme~$X$.
\end{corollary}

Taking~\Ref{V.2.3} into account, one can suppose $X$ \'etale over~$Y$. But if an $s \in G$ is in some $G_i(x)$, it then follows
from I~\Ref{I.5.4} that $s$ acts trivially on~$G$, hence is
the unit element since $G$ is faithful, which proves
the first assertion. Let $u$ be a $Y$-automorphism of~$X$, let
$x \in X$. By proposition~\Ref{V.1.3}, there exists an $s \in
G$ such that $s(x)=u(x)$, and inducing the same residue
homomorphism
$\kres(x)/\kres(x')$
as $u$. By \loccit one therefore has
$s=u$, which completes the proof.

\begin{remark}
\label{V.2.5}
The hypothesis that $G$ acts faithfully is obviously
not superfluous in corollary~\Ref{V.2.4}. The same holds
for the hypothesis that $X$ is connected, as one sees for instance by
taking $X=Y \times E$, $E$ being a finite set, and $G$ the
group of permutations of~$E$: $G$ acts with plenty of inertia,
nevertheless $(Y \times E)/G=Y \times (E/G)=Y$, and $X$ is \'etale
over~$Y$. Taking for $G$ a group strictly smaller than the
symmetric group of~$E$, but acting transitively on~$E$,
one sees that there will also be $Y$-automorphisms of~$X$ not coming
from~$G$.
\end{remark}

The typical example of a group $G$ acting without inertia is that of
$Y \times G$, on which one makes $G$ act by means of its
operations on the factor $G$ by right translations: a
$Y$-prescheme $X$ with right group of operators
$G$ is said to be \emph{trivial} if it is isomorphic to $Y \times G$.

To make the link between preschemes with finite groups
of operators and the notion of principal bundle in a
category (a link which, moreover, we shall not need for the
rest of the seminar, but which is important in other contexts) the
following considerations are useful. We fix a
base prescheme $Y$, and we place ourselves in the
category of $Y$-preschemes. If $G$ is a finite
% original p. 114
group, we shall set for brevity $G_Y=Y \times G$; it is thus a
finite group scheme over~$Y$ (\cf \numero 1), and if $X$ is
a $Y$-prescheme, one has
$$
X \times_Y G_Y = X \times G
$$
(same reference). The datum of a $Y$-morphism $X
\times_Y G_Y \to X$ is therefore equivalent to the datum of a
$Y$-morphism $X \times G \to X$, \ie to the datum, for every
$g \in G$, of a $Y$-morphism $T_g \colon X \to X$. One notes
immediately that for the datum of the $T_g$ to define on~$X$ a structure of prescheme with right group of operators
$G$ (\ie $T_{gg'}=T_{g'}T_g$, $T_e=\id_x$) it is necessary and
sufficient that the corresponding $Y$-morphism $X \times_Y G_Y \to X$
define on~$X$ a structure of~$Y$-prescheme with
$Y$-group scheme of operators (in the general sense
of objects with $\cal{C}$-group of operators in a
category $\cal{C}$). Suppose that this is so. Let us recall that
$X$ is said to be \emph{formally principal homogeneous}
under~$G_Y$\kern1pt\footnote{one now says rather: $X$ is a
pseudo-torsor under $G_Y$.} if the canonical morphism
$$
X \times_Y G_Y \to X \times_Y X
$$
whose components are respectively $\pr_1$ and the
multiplication morphism $\pi \colon X \times_Y G_Y \to X$, is an
isomorphism. In the present case, identifying the left-hand side with $X
\times G$, the morphism considered is the one which, to every $g
\in G$, associates the morphism
$$
(\id_X,T_g)=(\id_X \times_Y T_g)\Delta_{X/Y} \colon X \to X \times_Y X
$$
and consequently, to say that $X$ is formally principal homogeneous
under $G_Y$ also means that $X \times_Y X$ is isomorphic to the
direct sum of the transforms of the diagonal by the
elements $(e,g)$ of~$G \times G$ (acting on~$X \times_Y
X$ in the obvious way), where $e$ denotes the unit
element of~$G$. If one wishes not to distinguish left and right,
and to give a formula that remains applicable to the product of more than two
factors identical to $X$, one can formulate the condition by saying
that the canonical morphism
$$
X \times_G (G \times G) \to X \times_Y X
$$
obtained by attaching to the pair $(g,g')$ the morphism
$$
(T_g,T_{g'})=(T_g \times_Y T_{g'})\Delta_{X/Y} \colon X \to X \times_Y
X
$$
and by making $G$ act on the left on~$G \times G$ via
the diagonal homomorphism;
$$
s(g,g')=(sg,sg'),
$$
is an \emph{isomorphism}.

The
% original p. 115
notion of \emph{principal homogeneous space}
is deduced from that of formally principal
homogeneous space by adding a supplementary axiom, ensuring that the
``quotient'' of~$X$ by $G_Y$ exists and is precisely
the right unit object of the category, here $Y$. This
axiom may vary according to the context, and is often made explicit most
conveniently (in the yoga of ``descent'') by requiring that
the object with operators become ``trivial'', \ie isomorphic to the
product $X \times_Y G_Y$ (in the present case $X \times G$) by a
suitable change of base, of a specified type (in such a way, in practice, as to permit the technique of descent;
\cf Grothendieck, Technique de descente et th\'eor\`emes
d'existence en G\'eom\'etrie Alg\'ebrique, S\'em. Bourbaki
\No 190, pages 26 to 28)\footnote{\Cf Exp\ptbl \Ref{VIII} for the
theory of flat descent.}. In this order of ideas,
let us point out here the characterization of principal
homogeneous bundles with group $G$ (in the sense of \loccit):

\begin{proposition}
\label{V.2.6}
Let $Y$ be a locally noetherian prescheme, $X$ a
$Y$-prescheme with finite group $G$ of operators
acting on the right. The following conditions are
equivalent:
\begin{enumerate}
\item[(i)] $X$ is finite over~$Y$, $Y=X/G$, the inertia groups of the
points of~$X$ are reduced to the unit element.
\item[(ii)] There exists a faithfully flat and
quasi-compact change of base $Y_1 \to Y$ such that $X_1=X \times_Y Y_1$ is a
trivial $Y_1$-prescheme with operators, \ie isomorphic
to $Y_1 \times G$.
\item[(ii~bis)] Like (ii), but with $Y_1 \to Y$ being finite,
\'etale, surjective.
\item[(iii)] $X$ is formally principal homogeneous under $G_Y$,
and faithfully flat and quasi-compact over~$Y$.
\end{enumerate}
\end{proposition}
\textit{Proof} (i)$\To$(ii~bis) We shall take
$Y_1=X$, noting that $X \to Y$ is indeed finite, \'etale by~\Ref{V.2.3},
and surjective. Let us show that $X_1$ is then trivial over~$Y_1$, which
will follow from the

\begin{corollary}
\label{V.2.7} If \textup{(i)} is satisfied and if $X$ admits a
section over~$Y$, then $X$ is a trivial space with operators.
\end{corollary}

Indeed, this section makes it possible to define a $G$-morphism $X
\times G \to X$, surjective since $G$ is transitive on the fibers of
$X$, injective since $G$ acts without inertia; finally, it is a
local isomorphism by virtue of I~\Ref{I.5.3} since $X$ is \'etale
over~$Y$. Hence it is an isomorphism.

(ii~bis)
% original p. 116
trivially implies (ii), which implies (i) because the ingredients
of (i) are ``invariant'' under faithfully flat
quasi-compact extension of the base (\cf the S\'eminaire Bourbaki cited
above for ``finite''; for the inertia groups, one
applies~\Ref{V.2.1}, and for $Y=X/G$, a converse
of~\Ref{V.1.9} in the case of a \emph{faithfully flat} change of base,
which we had forgotten to make explicit).

We have proved (i)$\To$(iii) in passing, while proving
(i)$\To$(ii~bis). Finally (iii)$\To$(ii), because the
first hypothesis in (iii) means precisely that
$X$ becomes trivial on making the change of base $Y_1=X$; whence
(ii) since $X$ is faithfully flat and quasi-compact over~$Y$.

\begin{definition}
\label{V.2.8} A $Y$-prescheme $X$ with right group
of operators $G$ satisfying the
equivalent conditions of~\Ref{V.2.6} is called a \emph{principal
covering of~$Y$, with Galois group~$G$}.
\end{definition}

\section{Automorphisms and morphisms of \'etale coverings}
\label{V.3}

\begin{proposition}
\label{V.3.1} Let $X$ be \'etale, separated, of finite type over~$Y$, with $Y$ locally noetherian, and let $G$ be a finite group acting on~$X$ by $Y$-automorphisms. Then $G$ acts in an
admissible way and the quotient prescheme $X/G$ is \'etale over~$Y$.
\end{proposition}

One does not suppose $X$ finite over~$Y$; however $X$ is quasi-projective
over~$Y$, whence the existence of~$X/G$ thanks
to~\Ref{V.1.8}. Let us first prove

\begin{corollary}
\label{V.3.2} The morphism $X \to X/G$ is \'etale.
\end{corollary}

One can obviously suppose $G$ transitive on the set of
connected components of~$X$, then, by considering the
stabilizer of a connected component, that $X$ is even
connected. Finally, one can suppose that $G$ acts faithfully.
But then one sees as in~\Ref{V.2.4} that $G$ acts without
inertia, hence by~\Ref{V.2.3} it follows that $X \to X/G$ is
\'etale. One concludes thanks to the

\begin{lemma}[afterthought to Expos\'e I]
\label{V.3.3} Let $X \to X' \to
Y$ be morphisms of finite type, $x$ a point of~$X$, $x'$ and $y$ its
images. Suppose $Y$ locally noetherian. If two of the
morphisms considered are \'etale at the marked points, so is
the third.
\end{lemma}

It remains only to look at the case where $X \to X'$ and $X \to
Y$ are \'etale at $x$, and to prove that $X' \to Y$ is \'etale at $x'$ (which
is the case that we need
% original p. 117
for~\Ref{V.3.1}). Making a suitable flat extension of the base
$Y$, one is reduced to the case where the residue extension
$\kres(x)/\kres(y)$
is trivial. Consider the homomorphisms
$\cal{O}_y \to \cal{O}_{x'} \to \cal{O}_x$ and the homomorphisms
deduced by passing to the completions; the hypothesis
means that $\hat{\cal{O}_y} \to \hat{\cal{O}_x}$ and
$\hat{\cal{O}_{x'}} \to \hat{\cal{O}_x}$ are isomorphisms,
whence immediately that $\hat{\cal{O}_y} \to \hat{\cal{O}_{x'}}$
is one, which proves the lemma.

\begin{corollary}
\label{V.3.4} If $X$ is finite and \'etale over~$Y$, then $X/G$
is finite and \'etale over~$Y$.
\end{corollary}

\begin{proposition}
\label{V.3.5} Let $X$, $X'$ be two \'etale coverings of
$Y$. Then every $Y$-morphism $f \colon X \to X'$ factors as the
product of a surjective \'etale morphism $X \to X''$ and
the canonical immersion $X'' \to X'$ of a subset $X''$ of~$X'$ that is
both open and closed.
\end{proposition}

One knows (I~\Ref{I.4.8}) that $f$ is \'etale, hence an open
morphism; on the other hand, $X$ being finite over~$Y$, $f$ is closed,
hence $f(X)=X''$ is a subset of~$X'$ that is both open and closed. We are done (N.B. it would have sufficed that $X'$, instead of being an
\'etale covering, be unramified over~$Y$).

\begin{corollary}
\label{V.3.6} With the preceding notation, $X \to X'$
is a strict epimorphism in the category of
preschemes, and $X' \to X''$ is a monomorphism (and even
a strict monomorphism) in the category of preschemes.
\end{corollary}

The first assertion means by definition that the sequence of
morphisms
$$
\xymatrix@C=.7cm{X \times_{X''} X \ar@<2pt>[r]^-{\pr_1}
\ar@<-2pt>[r]_-{\pr_2} & X \ar[r] & X''}
$$
is exact, and this follows from the fact that $X \to X''$ is finite and
faithfully flat, as one sees easily (\cf Grothendieck, \loccit). The dual assertion for $X'' \to X'$ is even more trivial.

Corollary~\Ref{V.3.6} will be useful to us for the theory of the
fundamental group in the next no.; it is possible (for those
who do not like the notion of strict epimorphism) to replace
corollary~\Ref{V.3.6} by such a variant as the reader will arrange
to his personal taste. Let us merely take the opportunity to
point out that a factorization
% original p. 118
$f=f'f''$, with $f''$ a strict epimorphism and $f'$ a
monomorphism, is necessarily unique up to unique isomorphism
(in any category); however, there may exist at the
same time a factorization $f=f_1f_2$ having the
dual properties: $f_2$ is an epimorphism, $f_1$ a
strict monomorphism (also unique up to unique
isomorphism), which is not isomorphic to the preceding one: it
suffices to take for instance the category of topological vector spaces
(Hausdorff, if one insists), and for $u \colon X \to
X'$ a morphism such that $u(X)$ is not closed.

\begin{proposition}
\label{V.3.7} Let $Y$ be a \emph{connected} locally
noetherian prescheme, $y$ a point of~$Y$, $\mathit{\Omega}$ an
algebraically closed extension
of~$\kres(y)$. For every $X$ over~$Y$, one
denotes by $X(\Omega)$ the set of points of~$X$
with values in $\Omega$. Let $X$, $X'$ be \'etale coverings
of~$Y$ and $u \colon X \to X'$ a $Y$-morphism such that
the corresponding map $X(\Omega) \to
X'(\Omega)$ is an isomorphism. Then $u$ is an
isomorphism.
\end{proposition}

One is immediately reduced to the case where $X'$ is
connected. Since $X \to X'$ is finite and \'etale, one knows that the
geometric number of points in a fiber of~$X \to X'$ is
constant, and equal to 1 if and only if the morphism
considered is an isomorphism. Now this number is also the
number of elements in a fiber of~$X(\Omega) \to
X'(\Omega)$, whence the conclusion.

\section{Axiomatic conditions for a Galois theory}
\label{V.4}
Let $\cal{C}$ be a category, $F$ a covariant functor from
$\cal{C}$ into the category of finite sets. Suppose the
following conditions are satisfied:

\begin{enumerateb}

\item[(G 1)] $\cal{C}$ has a final object\footnote{Recall that an object
$e$ of~$\cal{C}$ is called a \emph{final object} if for every $X$ in
$\cal{C}$, $\Hom(X,e)$ has exactly one element. One defines
dually an \emph{initial object} of~$\cal{C}$.} and the fiber
product of two objects over a third one in $\cal{C}$
exists (this axiom can also be stated by saying that in
$\cal{C}$ finite projective limits exist).

\item[(G 2)] Finite sums exist in $\cal{C}$ (hence so does an
initial object $\emptyset_{\cal{C}}$, playing the role of the empty
set), as does the quotient of an object of~$\cal{C}$ by a finite
group of automorphisms.

\item[(G 3)] Let $u \colon X \to Y$ be a morphism in $\cal{C}$;
then $u$ factors as a product $X\lto{u'} Y'\lto{u''} Y$, with $u'$ a \emph{strict} epimorphism and
$u''$ a monomorphism which is an isomorphism onto a direct summand
of~$Y$.

\item[(G 4)]
% original p. 119
The functor $F$ is left exact (\ie transforms a right unit
into a right unit, and commutes with fiber
products).

\item[(G 5)] $F$ commutes with finite direct sums, transforms
strict epimorphisms into epimorphisms, and commutes with
passage to the quotient by a finite group of automorphisms.

\item[(G 6)] Let $u \colon X \to Y$ be a morphism in $\cal{C}$ such
that $F(u)$ is an isomorphism; then $u$ is an isomorphism.
\end{enumerateb}

Our object is to construct a topological group~$\pi$, a projective
limit of finite groups, and an equivalence of the
category $\cal{C}$ with the category $\cal{C}(\pi)$
\emph{of finite sets on which} $\pi$ \emph{operates
continuously} (\ie in such a way that the stabilizer of a
point is an open subgroup, or equivalently that there exists a discrete
quotient group which already operates on the set
under consideration), the equivalence constructed transforming the given
functor $F$ into the obvious inclusion functor of~$\cal{C}(\pi)$
into the category of finite sets. Let us note right away that
the category $\cal{C}(\pi)$ constructed by means of a topological
group~$\pi$, and the preceding inclusion functor,
do indeed satisfy conditions (G 1) to (G 6).

We proceed in several steps.
\begin{enumerateb}
\item[a)] \emph{Let} $u \colon X \to Y$ \emph{be in}
$\cal{C}$. \emph{For} $u$ \emph{to be a monomorphism, it is necessary and
sufficient that} $F(u)$ \emph{be one}. (Uses (G 1), (G 4), (G 6)).

Indeed, to say that $u$ is a monomorphism means that the projection
$\pr_1 \colon X \times_Y X \to X$ is an isomorphism.

\item[b)] \emph{Every object} $X$ \emph{of} $\cal{C}$ \emph{is
artinian}.

Indeed, if $X' \to X'' \to X$ are monomorphisms such that
$F(X')$ and $F(X'')$ have the same image in $F(X)$, then by a)
$F(X') \to F(X'')$ is an isomorphism, hence $X' \to X''$ is an
isomorphism by (G 6).

\item[c)] \emph{The functor} $F$ \emph{is strictly
pro-representable}. (\cf Grothendieck, Technique de descente et
th\'eor\`emes d'existence en G\'eom\'etrie Alg\'ebrique, II,
S\'eminaire Bourbaki 195, February 1960).

Indeed, by \loccit prop\ptbl \Ref{V.3.1}, this follows from
b) and (G 4). One can therefore find a projective system over a
filtered ordered set~$I$:
$$
P=(P_i)_{i \in I}
$$
% original p. 120
in $\cal{C}$, considered as a pro-object of~$\cal{C}$, and a
functorial isomorphism
\begin{equation*}
\label{eq:V.4.*} \tag{$*$} F(X)=\Hom_{\Pro(\cal{C})}(P,X)=\varinjlim_{i} \Hom_{\cal{C}}(P_i,X)
\end{equation*}
This isomorphism is realized by an element
$$
\varphi \in \varprojlim_{i} F(P_i)=F(P)
$$
One may moreover assume that the transition homomorphisms
$\varphi_{ji} \colon P_i \to P_j$ ($i \geq j$) are
\emph{epimorphisms}, and that \emph{every epimorphism} $P_i \to
P'$ is equivalent to an epimorphism $P_i \to P_j$ for
suitable $j \leq i$ (which determines the projective system~$P$ in an essentially unique way).

An object $X \in \cal{C}$ is said to be \emph{connected}
if it is not isomorphic to the sum of two objects of~$\cal{C}$
not isomorphic to the initial object $\emptyset_{\cal{C}}$.

\item[d)] \emph{The} $P_i$ \emph{are connected and not isomorphic
to} $\emptyset_{\cal{C}}$.

If $X$ is a left unit, one has $F(X)=\emptyset$ by (G 5)
applied to the direct sum of an empty family, and
conversely by (G 6). Hence if $X'$ is an object of~$\cal{C}$
which is not a left unit, \ie such that $F(X') \neq
\emptyset$, there exists no morphism from~$X'$ into $X$. Hence if some
$P_i$ is a left unit, then $i$ is a largest
element of the filtered ordered index set $I$, and
formula \eqref{eq:V.4.*} would mean $F(X)=\Hom(P_i,X)={}$\kern2pt a set
reduced to one element for every $X$, which is absurd
since $F(\emptyset_{\cal{C}})=\emptyset$. Hence the~$P_i$ are not
isomorphic to $\emptyset_{\cal{C}}$.

Suppose that one has $P_i=A \amalg B$, whence by (G 5)
$F(P_i)=F(A) \amalg F(B)$; in particular the element $a_i$
of~$F(P_i)$, corresponding via \eqref{eq:V.4.*} to
the identity homomorphism $P_i \to P_i$, lies in $F(A) \amalg
F(B)$, for instance in $F(A)$. This means that there exists a $j
\geq i$ such that $\varphi_{ij} \colon P_j \to P_i$ factors as
$P_j \to A \to P_i=A \amalg B$, where the second arrow
is the canonical morphism. Hence $F(P_j) \to F(P_i)$ factors
as $F(P_j) \to F(A) \to F(P_i)=F(A) \amalg F(B)$, and since $F(P_j)
\to F(P_i)$ is surjective by (G 5), it follows that
$F(B)=\emptyset$, hence $B$ is isomorphic to $\emptyset_{\cal{C}}$.

\item[e)]
% original p. 121
\emph{Every morphism} $u \colon X \to Y$ \emph{in} $\cal{C}$,
\emph{with} $X$ \emph{not isomorphic to} $\emptyset_{\cal{C}}$
\emph{and} $Y$ \emph{connected, is a strict epimorphism. Every
endomorphism of a connected object is an automorphism}.

Consider the factorization (G 3) of~$u$; since $X \neq
\emptyset_{\cal{C}}$ it follows from (G 6) that $F(X) \neq \emptyset$, hence
$F(Y') \neq \emptyset$, hence $Y' \neq \emptyset_{\cal{C}}$, hence, $Y$
being connected, $Y'$ is identified with~$Y$, hence $u$ is a
strict epimorphism. Suppose that $u$ is an endomorphism of
the connected object $X$; let us prove that it is an automorphism. Indeed, we
may assume~$X$ not isomorphic to $\emptyset_{\cal{C}}$, hence $u$ is
a strict epimorphism by the foregoing, hence $F(u)$ is
an epimorphism by (G 5), and since $F(X)$ is a finite set,
$F(u)$ is bijective. Hence $u$ is an automorphism by (G 6).

In particular, \emph{every endomorphism of a} $P_i$ \emph{is an
automorphism}.

\item[f)] \emph{The following conditions on a} $P_i$ \emph{are
equivalent}: (i) The natural injective map
$\Hom(P_i,P_i) \to \Hom(P,P_i) \simeq F(P_i)$ is also surjective,
\ie for every $u \colon P \to P_i$ there exists a $v \colon P_i \to
P_i$ such that $u=v\varphi_i$ (where $\varphi_i$ is the canonical
homomorphism $P \to P_i$). (ii) The group $\Aut(P_i)$ operates transitively on~$F(P_i)$. (iii)~The group $\Aut(P_i)$ operates
simply transitively on~$F(P_i)$.

Indeed, identifying $\Hom(P,P_i)$ with $F(P_i)$, the map
considered in (i) is none other than $v \mto
F(v)(\varphi_i)$. The equivalence of the three conditions then follows
from the fact that $\Hom(P_i,P_i)=\Aut(P_i)$ and that the preceding
map is already injective.

A $P_i$ satisfying the equivalent conditions
(i), (ii), (iii)
of~f) is called \emph{Galois}.

\item[g)] \emph{For every} $X$ \emph{in} $\cal{C}$, \emph{there exists
a} Galois $P_i$ \emph{such that every} $u \in \Hom(P,X)$ \emph{factors
as} $P \lto{\varphi_i} P_i \longrightarrow X$.

Let $J=\Hom(P,X)=F(X)$; it is a finite set, hence there exists a
$P_j$ such that every $u \colon P \to X$ factors as $P\lto{j}P_j\longrightarrow X$, or equivalently such that the natural morphism
$$
P \to X^J \qquad (J=\Hom(P,X))
$$
factors as
$$
P \lto{\varphi_j} P_j \to X^J
$$
By
% original p. 122
virtue of (G~3), the morphism $P_j \to X^J$ factors as the product
of a monomorphism and a strict epimorphism, which one may
take of the form $\varphi_{ij} \colon P_j \to P_i$. We are therefore
reduced to proving that $P_i$ is Galois. Let $k$ be an index
$\geq j$ such that every morphism $P \to P_i$ factors through
$P \lto{\varphi_k} P_k \longrightarrow P_i$. Note that the
natural morphism $P_k \to X^J$ again factors as the composite
$$
P_k \lto{\varphi_{ik}} P_i \lto{U} X^J
$$
where the first arrow is a strict epimorphism
by~e), and the second a monomorphism. We want to prove that for
a given morphism $\psi \colon P_k \to P_i$, there exists an
endomorphism $v$ of~$P_i$ such that $\psi = v \varphi_{ik}$. But for
every $u \in \Hom(P_i,X)$, consider $u \psi \in \Hom(P_k,X)$;
it is therefore of the form $u'\varphi_{ik}$ with a well-determined
$u' \in \Hom(P_i,X)$. The map $u \mto u'$ from~$J$ into $J$
thus defined by $\psi$ is moreover injective, since $\psi$ is
an epimorphism by virtue of e); it is therefore bijective since
the set~$J$ is finite. The bijective map $u \mto u'$ from~$J$
into $J$ therefore defines an isomorphism $\alpha \colon X^J \isomto
X^J$ making the diagram
$$
\xymatrix{P_k \ar[r]^-{\varphi_{ik}} \ar@{=}[d] & P_i \ar[r]^-U & X^J
\ar[d]^-{\alpha}_-{\simeq} \\ P_k \ar[r]^{\psi} & P_i \ar[r]^-U & X^J}
$$
commutative. By the uniqueness properties of the factorization
of a morphism as a product of a monomorphism and a strict
epimorphism, it follows (since $\psi$ too is a strict
epimorphism by e)) that one can find a morphism $v \colon P_i \to
P_i$ which leaves the diagram commutative, qed.

We conclude in particular that \emph{the Galois} $P_i$
\emph{form a cofinal system in the system of the} $(P_j)$. We
shall therefore have, since for a Galois object $P_i$ one has
$$
\Hom(P,P_i)=\Hom(P_i,P_i)=\Aut(P_i),
$$
by passing to the limit:
$$
\Hom(P,P) = \varprojlim_{i} \Hom(P,P_i) = \varprojlim_{i} \Hom(P_i,P_i) =
\varprojlim_{i} \Aut(P_i)
$$
where
% original p. 123
the projective limit is taken over the Galois $P_i$. Moreover,
by means of the identification $\Hom(P,P_i)=F(P_i)$, and taking into account that
$F$ transforms epimorphisms into epimorphisms, one sees that the
transition homomorphisms in the preceding projective system
are surjective. We conclude from all this:

\item[h)] \emph{One has}
$$
\Hom(P,P) = \Aut(P) = \varprojlim_{i} F(P_i) = \varprojlim_{i} \Aut(P_i),
$$
\emph{where the projective limit is taken over the Galois}
$P_i$.

In particular, $\Aut(P)$ appears as the projective limit of a
projective system of finite groups (the transition homomorphisms
being surjective); we shall endow it with the projective limit
topology of the discrete topologies. \emph{We shall denote
by}~$\pi$
\emph{and call the fundamental group}
(of~$\cal{C}$ equipped with~$F$) \emph{the group opposite
to}~$\Aut(P)$. This group therefore operates \emph{on the right}
on~$P$; it is the projective limit of finite groups $\pi_i$
operating on the right on the Galois $P_i$, $\pi_i$ being
the group opposite to $\Aut(P_i)$.

Taking into account the functorial isomorphism
$$
F(X)=\Hom(P,X)
$$
and the definition of~$\pi$, one therefore sees that $\pi$ operates
\emph{on the left} on~$F(X)$, and moreover continuously,
by g) (for with the notation of g), it is in fact $\pi_i$
which operates on~$F(X)$). It is trivial that for every morphism $u
\colon X \to Y$ in $\cal{C}$, the morphism $F(u) \colon F(X) \to
F(Y)$ is compatible with the operations of~$\pi$. \emph{One may
therefore henceforth regard} $F$ \emph{as a covariant
functor}
$$
F \colon \cal{C} \to \cal{C}(\pi)
$$
\emph{where} $\cal{C}(\pi)$ \emph{is the category of finite
sets on which} $\pi$ \emph{operates continuously on the left}.

We shall now define a functor in the opposite direction:
$$
G \colon \cal{C} \from \cal{C}(\pi)
$$
by the formula
$$
G(E)=P \times_{\pi} E,
$$
where $P \times_{\pi} E$ is defined as the solution of the
universal problem summarized by
$$
\Hom_{\cal{C}}(P \times_{\pi} E,X) \isomto \Hom_{\pi}(E,\Hom(P,X))
$$
(where
% original p. 124
in the right-hand side $\Hom(P,X)=F(X)$ is regarded
as a set on which $\pi$ operates on the left). One must prove
the existence of the object $P \times_{\pi} E$.

\item[i)] \emph{Let} $Q$ \emph{be an object of} $\cal{C}$ \emph{on which a
finite group} $G$ \emph{operates on the right, and} $E$ \emph{a
finite set on which} $G$ \emph{operates on the left. Then} $G
\times_G E$ \emph{exists, and the canonical map}
$$
F(Q) \times_G E \to F(Q \times_G E)
$$
\emph{is an isomorphism}.

Since finite direct sums exist in $\cal{C}$ by (G 2), and
$F$ commutes with them by (G 5), one is immediately reduced to the case
where $G$ operates transitively on~$E$, for if the~$E_j$ are the
orbits of~$G$ in $E$, one will have
$$
Q \times_G E = \underset {j}{\amalg}\ Q \times_G E_j.
$$
Let then $a \in E$, and let $H$ be its stabilizer; one sees immediately
from the definition that $Q \times_G E$ is identified with $Q/H$.
Whence the existence, thanks to (G~2), and the commutation property
for~$F$, thanks to (G~5).

\item[j)] \emph{Let} $E$ \emph{be an object of} $\cal{C}(\pi)$, \emph{and
let} $P_i$ \emph{be Galois such that} $\pi_i$ \emph{already
operates on} $E$. \emph{Then} $P_i \times_{\pi_i} E$
\emph{exists and one has a canonical isomorphism}
$$
E \isomto F(P_i \times_{\pi_i} E)
$$
\emph{If} $j \geq i$ \emph{is such that} $P_j$ \emph{is Galois,
then the canonical homomorphism} $P_j \times_{\pi_j} E \to\allowbreak P_i
\times_{\pi_i} E$ \emph{is an isomorphism}.

The first assertion is a special case of i), taking into account
that $\pi_i$ operates simply transitively on~$F(P_i)$,
which is equipped with a marked point $\varphi_i$, whence an isomorphism
$F(P_i) \times_{\pi_i} E \simeq E$. For the second assertion one
uses for instance (G 6).

For every $j$, let $\cal{C}_j$ be the full subcategory of
$\cal{C}$ formed by the $X$ such that $\Hom(P_j,X) \to \Hom(P,X)
\simeq F(X)$ is bijective. We know by g) that $\cal{C}$ \emph{is
the filtered union of the} $\cal{C}_j$. We therefore have, for $X \in
\cal{C}_j$:
\begin{align*}
\Hom_{\pi}(E,\Hom(P,X)) & \simeq \Hom_{\pi}(E,\Hom(P_j,X)) \simeq
\Hom_{\pi_j}(E,\Hom(P_j,X))\\ & \simeq \Hom(P_j \times_{\pi_j} E,X)
\end{align*}
and,
% original p. 125
taking into account the last assertion in j), we find an
isomorphism functorial in the object $X$ of~$\cal{C}_j$:
$$
\Hom_{\pi}(E,\Hom(P,X)) \simeq \Hom(P_i \times_{\pi_i} E,X)
$$
Since this is true for every $j$ and these functorial isomorphisms,
for variable $j$, induce one another, we conclude:

\item[k)] \emph{Under the conditions of} j), \emph{the composite
of the canonical morphisms}
$$
E \to \Hom(P_i,P_i \times_{\pi_i} E) \to \Hom(P,P_i \times_{\pi_i} E)
$$
\emph{makes} $P_i \times_{\pi_i} E$ \emph{a solution of the
universal problem defining} $P \times_{\pi} E$,
\ie \emph{the latter exists and one has an isomorphism}
$$
P \times_{\pi} E \isomto P_i \times_{\pi_i} E
$$

This completes the construction of the functor $G(E)$. On the other hand, one has
a functorial homomorphism
$$
\alpha \colon \id_{\cal{C}(\pi)} \to FG
$$
\ie a homomorphism functorial in the object $E$ of~$\cal{C}(\pi)$:
$$
\alpha(E) \colon E \to FG(E) = F(P \times_{\pi} E)
$$
namely the composite of the canonical morphisms
$$
E \to F(P) \times_{\pi} E \to F(P \times_{\pi} E)
$$
(where the first comes from the marked point $\varphi \in F(P)$).
Combining j) and k), one finds:

\item[l)] \emph{The homomorphism} $\alpha$ \emph{is an isomorphism}

One defines likewise a functorial homomorphism
$$
\beta \colon GF \to \id_{\cal{C}}
$$
\ie a homomorphism functorial in the object $X$ of~$\cal{C}$:
$$
\beta(X) \colon P \times_{\pi} F(X) \to X
$$
as
% original p. 126
the one associated with the $\pi$-homomorphism
$$
F(X) \to \Hom(P,X)
$$
inverse to the canonical isomorphism $\Hom(P,X) \isomto F(X)$.

\item[m)] \emph{The composites}
\begin{gather*}
F(X) \lto{\alpha(F(X))} FGF(X) \lto{F(\beta(X))} F(X)\\
G(E) \lto{G(\alpha(E))} GFG(E) \lto{\beta(G(E))} G(E)
\end{gather*}
\emph{are the identity isomorphisms}.

A trotting donkey.

Taking l) into account, it follows:

\item[n)] \emph{The homomorphism} $\beta$ \emph{is an isomorphism}

We have thus obtained the promised result:
\end{enumerateb}

\begin{theorem}
\label{V.4.1}
Let $\cal{C}$ be a category satisfying conditions \emph{(G
1)}, \emph{(G 2)}, \emph{(G~3)} at the beginning of this no., and $F$ a
covariant functor from~$\cal{C}$ into the category of finite
sets, satisfying conditions \emph{(G 4)}, \emph{(G 5)} and
\emph{(G 6)}. Then the preceding canonical constructions
define equivalences of categories $F\colon
\cal{C}\to \cal{C}(\pi)$ and $G\colon \cal{C}(\pi)\to \cal{C}$
quasi-inverse to one another. More precisely, there exists
a pro-object~$P$ of~$\cal{C}$, and a functorial isomorphism $F(X)
\isomfrom \Hom(P, X)$; $\pi$ is the group opposite to the group of
automorphisms of~$P$, topologized in a suitable way,
so that~$\pi$ operates continuously on the
sets $\Hom(P, X)\simeq F(X)$. Finally $G$ is given by
$G(E)\simeq P\times_{\pi} E$.
\end{theorem}

\begin{remarks}
\label{V.4.2}
The statement of conditions (G 1) to (G 6) becomes simpler
and more pleasant if one replaces (G 2) and (G 5) respectively by:
\begin{enumerateb}
\item[(G' 2)] Finite inductive limits exist in
$\cal{C}$.
\item[(G' 5)] The functor $F$ is right exact (\ie commutes with
finite inductive limits).
\end{enumerateb}
These conditions are apparently stronger than (G~2) and (G~5), but
it follows immediately from the structure
theorem~\Ref{V.4.1} that they are implied by (G~1)
to (G~6). Note, however, that in the cases which will
interest us, the verification of (G~2) and (G~5) seems
in fact simpler than that of (G'~2) and (G'~5). I do not know whether,
in condition (G~3), the fact that $u''$ is an isomorphism onto a
direct summand of~$Y$ could not be omitted.
\end{remarks}

\section{Galois categories}
\label{V.5}
% original p. 127

\begin{definition}
\label{V.5.1}
A Galois category is
a category $\cal{C}$ equivalent to a category
$\cal{C}(\pi)$, where $\pi$ is a compact group, projective limit
of finite groups (\ie totally disconnected).

For the definition of~$\cal{C}(\pi)$, \cf the beginning of \No \Ref{V.4}. By
Theorem~\Ref{V.4.1}, $\cal{C}$ is Galois if and only if
it satisfies conditions (G 1) to (G 3), and if there exists a
functor $F$ from
$\cal{C}$
into the category of finite sets satisfying conditions
(G~4) to (G~6) (\ie which is \emph{exact} and \emph{conservative},
in a general terminology). Such a functor will be
called a \emph{fundamental functor}
of the Galois category $\cal{C}$\kern1pt\footnote{It seems preferable to adopt the more suggestive term
``fiber functor''.};
it is pro-representable by a pro-object which we shall denote by~$P_F$;
a pro-object $P$ such that the associated functor $F$ is fundamental
is called a \emph{fundamental pro-object}.
In this way, the category of fundamental functors on~$\cal{C}$ is anti-equivalent to the category of
fundamental pro-objects; if $F$ and $P$ correspond to each other, the group
$\Aut F$ is therefore isomorphic to the opposite of the group $\Aut P$,
hence the group denoted $\pi$ in the preceding no.\ is
none other than $\Aut P$. Let us recall that in the preceding no.\ we
constructed, starting from a \emph{given} fundamental functor
$F$, an equivalence of~$\cal{C}$ with
$\cal{C}(\pi)$ (where $\pi=\Aut(F)$) which transforms the functor $F$
into the canonical functor from~$\cal{C}(\pi)_x$ into the category of
finite sets. In this typical case $\cal{C}=\cal{C}(\pi)$, $F=$
canonical functor, the fundamental pro-object associated with $F$
is none other than the projective system of the discrete quotients
$\pi_i$ of~$\pi$.
\end{definition}

It may be useful to make explicit the category of pro-objects of
$\cal{C}(\pi)$. One finds:

\begin{proposition}
\label{V.5.2}
The category $\Prodash\cal{C}(\pi)$
is canonically equivalent to the category
$\cal{C}'(\pi)$ of spaces, with topological group $\pi$
of operators, which are compact and totally disconnected.
\end{proposition}

Since the latter contains $\cal{C}(\pi)$ as a
full subcategory (corresponding to the compact discrete spaces with
operators) and since projective limits
exist in it, one has in any case a canonical functor $g\colon
\Prodash\cal{C}(\pi) \to \cal{C}'(\pi)$, associating with the projective system
$Q=(Q_i)$
the object $X=\varprojlim_{i} Q_i$ of~$\cal{C}'(\pi)$. To define
a functor in the opposite direction, it suffices to define a contravariant
functor from~$\cal{C}'(\pi)$ into the category of left exact functors
$\cal{C}\to \Ens$, and we shall take for every
$X\in\cal{C}'(\pi)$ the functor $h(X)(E)= \Hom(X, E)$
% original p. 128
(the $\Hom$ being taken in $\cal{C}'(\pi)$). It is immediate
by definition that the functors $h$ and $g$ are adjoint to each
other, and that $hg$ is canonically isomorphic to the identity functor
of~$\Prodash\cal{C}(\pi)$. It remains to prove (in order to establish
that $g$ and $h$ are quasi-inverse to each other) that every object of
$\cal{C}'(\pi)$ is isomorphic to an object of the form $g(Q)$,
where $Q\in\Prodash{\cal C}(\pi)$, in other words: \emph{every
space $X$ with topological group $\pi$ of operators, which is
compact and totally disconnected, is isomorphic to a projective
limit of finite discrete spaces with operators}. Since $X$
is the projective limit of its finite discrete quotients (as a
topological space without operators), we are reduced to
showing that in the set of these quotients, there is a
cofinal system which is invariant under $\pi$. For this it suffices to show
that for such a quotient $X'$, the set of transforms of this
quotient by the operations of~$\pi$ is finite, (one will then take
the sup of the said transforms, which will be an invariant quotient
majorizing $X'$). Or again, that there is an open invariant subgroup
$\pi'$ of~$\pi$ such that the elements of this subgroup
leave $X'$ invariant. Now $X'$ corresponds to a finite partition of~$X$ into
open sets $X_i$. By reason of continuity and of
compactness of~$\pi$, there exists a neighborhood $V$ of
the neutral element of~$\pi$ such that $s\in V$ implies
$s\cdot X_i\subset X_i$ for every~$i$, hence $s$ leaves~$X'$ invariant. Now we know
that the open invariant subgroups of~$\pi$ form a
fundamental system of neighborhoods of the neutral element, which
completes the proof.

Let us note that one sees even more simply that the category
$\Ind \cal{C}(\pi)$ is canonically equivalent to the
category of sets on which $\pi$ operates
continuously. We shall not need this here.

\begin{proposition}
\label{V.5.3}
Let $\cal{C}$ be a Galois category, $F$ a fundamental
functor
on~$\cal{C}$,
$P=(P_i)$ the associated pro-object, normalized in the usual
way. Let
$X\in\cal{C}$;
for $X$ to be connected, it is necessary and sufficient that $\pi$ operate
transitively on~$E=F(X)$.
\end{proposition}

We are reduced to the typical case $\cal{C}=\cal{C}(\pi)$, $F=$ canonical
functor, where it is trivial.

\begin{corollary}
\label{V.5.4}
Equivalent conditions on~$X$: \textup{(i)} $X$ is connected and
$\not\simeq\emptyset_{\cal{C}}$ \textup{(ii)} the group $\pi$ is
transitive on~$E=F(X)$, and $F(X)\ne\emptyset$ \textup{(iii)} $X$ is
isomorphic to a $P_i$.
\end{corollary}

The equivalence
% original p. 129
of (i) and (iii) also already follows easily from
\No \Ref{V.4},~e).

\begin{proposition}
\label{V.5.5}
Let $Q=(Q_i)_{i\in I}$ be a pro-object
of~$\cal{C}$,
normalized in the usual way, and let $G$ be the
corresponding functor $G(X)=\Hom(Q, X)$ from~$\cal{C}$
(into~$\Ens$).
The following conditions are equivalent:
\begin{enumerate}
\item[(i)] $G$ commutes with finite direct sums
\item[(ii)] $G$ commutes with the sum of two objects
\item[(iii)] The $Q_i$ are connected and
$\not\simeq\emptyset_\cal{C}$
\item[(iv)] $Q$ is isomorphic to $\pi/H$, where $H$ is a
closed subgroup of~$\pi$.
\item[(v)] The functor $G$ is isomorphic to the functor $E\mto
E^H$ (set of invariants under $H$) defined by a closed
subgroup $H$ of~$\pi$.
\end{enumerate}
\end{proposition}

N.B. in the statement of (iv) and (v), one assumes chosen a
fundamental functor, making it possible to identify $\cal{C}$ with the
category~$\cal{C}(\pi)$.

\subsubsection*{Proof} We may assume
$\cal{C}=\cal{C}(\pi)$. The implication (i)$\To$(ii) is
trivial, (ii)$\To$(iii) is proved like
property d) of \No \Ref{V.4}. Let us prove (iii)$\To$(iv). Indeed,
we know that $\varprojlim_i Q_i$ is nonempty as a projective
limit of nonempty finite sets. Let $a$ be a point of
$\varprojlim_i Q_i$; it defines a homomorphism of spaces with
operators
$$
\pi\to Q
$$
which is \emph{surjective}, since for every $i$ the composite $\pi\to
Q\to Q_i$ is, because $\pi$ is transitive on $Q_i$ by virtue
of~\Ref{V.5.3}. If therefore $H$ is the stabilizer subgroup of
$a$, one obtains an isomorphism $\pi/H\isomto Q$. The implications
(iv)$\To$(v) and (v)$\To$(i) are again
trivial.

\begin{proposition}
\label{V.5.6}
Let $\cal{C}$ be a Galois category, $P$ a fundamental
pro-object of~$\cal{C}$ and $F$ the associated fundamental
functor. Let $P'=(P'_i)_{i\in I}$ be a pro-object of~$\cal{C}$, put
in normal form, and $F'$ the associated functor $F'(X)=\Hom(P',X)$
from~$\cal{C}$ into~$\Ens$. Equivalent conditions:
\begin{enumerate}
\item[(i)] $P'\simeq P$, or again $F'\simeq F$.
\item[(ii)] $P'$ is fundamental, or again $F'$ is
fundamental.
\item[(iii)] $F'$ transforms a sum of two objects into a sum, and
$X\not\simeq\emptyset_{\cal{C}}$ implies $F(X)\ne\emptyset$.
\item[(iv)]
% original p. 130
The objects of~$\cal{C}$ that are connected and $\ne\emptyset_{\cal{C}}$ are
exactly the objects isomorphic to a~$P'_i$.
\end{enumerate}
\end{proposition}

We have trivially (i)$\To$(iii) and (i)$\To$(ii),
moreover (ii)$\To$(iv) by virtue of~\Ref{V.5.4} (applied
to $P'$ instead of~$P$). Moreover (iii) or (iv) implies by virtue
of~\Ref{V.5.5} that $P'$ is of the form $\pi/H$ where $H$ is a
closed subgroup of~$\pi$. In case (iii), there exists for every
open invariant subgroup $\pi'$ of~$\pi$ a $\pi$-homomorphism
$P'=\pi/H\to \pi/\pi'$, whence $H\subset \pi'$, whence $H=(0)$ and
consequently (i), qed.

\begin{corollary}
\label{V.5.7}
Let $\cal{C}$ be a Galois category. The fundamental
pro-objects are isomorphic, the fundamental functors are
isomorphic.
\end{corollary}

In other words, \emph{the category of fundamental functors
is a connected groupoid~$\Gamma$},
which one may call the \emph{fundamental groupoid}
of the Galois category~$\cal{C}$. If $\cal{C}=\cal{C}(\pi)$,
the group of automorphisms of an object of the fundamental groupoid
is isomorphic to $\pi$, this isomorphism being well
determined up to inner automorphism. (N.B. one
calls a \emph{groupoid} a category in which all the
morphisms are isomorphisms, a \emph{connected} groupoid a
groupoid all of whose objects are isomorphic). The fundamental
pro-objects of~$\cal{C}$ form a connected groupoid
equivalent to the \emph{opposite} of the fundamental
groupoid. If $F, F'$ are two fundamental functors, associated
with fundamental pro-objects $P, P'$, then $\Hom(F, F')=\Isom(F,
F')$ is sometimes denoted $\pi_{F', F}$ and plays the role of a
``set of \emph{classes of paths} from~$F$ to~$F'$''. In
particular $\pi_{F, F}=\pi_F$ is none other than the \emph{fundamental
group of~$\cal{C}$} at $F$ constructed in the preceding
no. As for the pro-object $P$ associated with~$F$, it
plays the role of a \emph{universal covering at~$F$} of
the final object $e_{\cal{C}}$ of~$\cal{C}$.

It may be convenient to have a description of~$\cal{C}$ (up to
equivalence) in terms of its fundamental groupoid
$\Gamma$, without going through a choice of a particular object $F$ of the
latter. Now to every object $X$ of~$\cal{C}$ is associated the
functor $E_X$ on the fundamental groupoid, defined by
$$
E_X(F)=F(X),
$$
with values in~$\Ens$. (Such a functor is known in topology
under the name of ``local system'' on the groupoid)
$F(X)=E_X(F)$ may be called the \emph{fiber} of~$X$ at $F$,
and the functor $E_X$ the fiber-functor associated with~$X$. The
functor $E_X$ has the following
% original p. 131
property:
\emph{for every $F$, $E_X(F)$ is a finite set
on which the topological group $\pi_F=\Aut(F)$ operates continuously}.
For a
given covariant functor $\xi$ from the fundamental groupoid
into~$\Ens$, the preceding condition is moreover equivalent
to the same condition for \emph{one} arbitrary fixed $F$.
This being so:

\begin{proposition}
\label{V.5.8}
The functor $X\mto E_X$ is an equivalence of the category
$\cal{C}$ with the category of covariant functors from the
fundamental groupoid $\Gamma$ of~$\cal{C}$ into~$\Ens$, which
satisfy the condition
highlighted above.
\end{proposition}

Indeed, let $F_0$ be an object of the fundamental groupoid, and let
$\pi_0=\Aut(F_0)$; then the functor $\xi\mto \xi(F_0)$ is an
equivalence of the second category considered
in~\Ref{V.5.8} with the category $\cal{C}(\pi_0)$, as one
sees at once. On the other hand, the composite of the latter and
of the functor $X\mto E_X$ is the natural equivalence
$\cal{C}\to\cal{C}(\pi_0)$. It follows that the functor
$X\mto E_X$ itself is an equivalence.

\begin{corollary}
\label{V.5.9}
The category $\Prodash\cal{C}$
is canonically equivalent to the category of covariant
functors $\xi$ from the fundamental groupoid $\Gamma$ into the
category of topological spaces, satisfying the condition:
for every object $F$ of~$\Gamma$, $\xi(F)$ is a compact totally
disconnected space with topological group $\pi_F$
of operators.
\end{corollary}

Here again, one can verify this condition on~$\xi$; it suffices
to verify it for \emph{one}~$F$. The proof is the
same as for~\Ref{V.5.8}.

\begin{remark}
\label{V.5.10}
Let $(F_s)_{s\in S}$ be a family
of objects
of the fundamental groupoid~$\Gamma$. Let us set, for $s, s'\in S$:
$$
\Hom(s, s')=\Hom(F_s, F_{s'})
$$
so that $S$ itself becomes a connected groupoid and
the map $s\mto F_s$ a fully faithful functor from~$S$
into $\Gamma$, call it~$f$. Considering then the functor $X\mto
E_X\circ f$ from~$\cal{C}$ into the functor category
$\Hom(S,\Ens)$, one obtains a variant of~\Ref{V.5.8}
(and~\Ref{V.5.9}) with $\Gamma$ replaced by~$S$. The statement
thus obtained reduces to Theorem~\Ref{V.4.1} when $S$
is reduced to a point, and is none other than~\Ref{V.5.8}
itself if $S$ is the set of objects of~$\Gamma$.
\end{remark}

We are going to use~\Ref{V.5.9} to define a canonical
pro-object of~$\cal{C}$. For this, we consider the functor
from~$\Gamma$ into the category of topological spaces
% original p. 132
(and even of topological groups):
$$
f\colon F\mto \Aut(F)=\pi_F.
$$
This functor satisfies the condition considered in~\Ref{V.5.8};
the space with operators $f(F)$ under $\pi_F$ is none other than
$\pi_F$, considered as a space with operators under
itself by inner automorphisms. Hence the functor
$f$ corresponds to a pro-object of~$\cal{C}$ determined up to
unique isomorphism, which is even a pro-group of
$\cal{C}$ and which is called the \emph{fundamental pro-group
of~${\cal C}$},
playing the role of a local system of fundamental
groups. It is therefore a pro-group $\Pi$ of~$\cal{C}$ defined
by the condition that one have an isomorphism functorial in $F$
$$
F(\Pi)\simeq\pi_F.
$$

If $X$ is an arbitrary pro-object of~$\cal{C}$, one has a canonical
morphism
$$
\Pi\times X \to X
$$
which makes~$X$ an object with a group of operators on the left
$G$ in $\Prodash\cal{C}$. For this it suffices to note that for $F$
variable, one has a canonical map
$$
\Pi(F)\times X(F) \to X(F)
$$
\ie $$
\Aut(F)\times E_X(F) \to E_X(F), \quad \text{or} \quad \pi_F\times
F(X)\to F(X)
$$
which is functorial in~$F$. It is also functorial in $X$, hence
for every morphism $X\to Y$ of pro-objects, the corresponding
diagram
$$
\xymatrix{
\Pi \times X \ar[r]\ar[d]& X\ar[d] \\ \Pi \times Y \ar[r]& Y
}
$$
is commutative.

\begin{remark}
\label{V.5.11}
One must be careful not to confuse a fundamental pro-object $P$ (which is not
endowed with a group structure, and is connected) with the fundamental
pro-group (which is a pro-\emph{group}, and in general not
connected). Precisely, $G$ is connected if
and
only if $\pi_F$ operating on itself by inner
automorphisms is transitive, \ie if $\pi$ is reduced to
the neutral element, or again $\cal{C}$ equivalent to the
category of finite sets. Another
% original p. 133
essential difference is that $G$ is determined up to
unique isomorphism, and $P$ is determined only up to
non-unique isomorphism.

Let $E$ be a finite set and let us consider the constant functor on
the groupoid~$\Gamma$ with value $E$: by virtue
of~\Ref{V.5.8} it defines an object of~$\cal{C}$, denoted $E_{\cal{C}}$, which
may also be interpreted as the sum of~$E$ copies of
the final object $e_{\cal{C}}$ of~$\cal{C}$. One may consider
$E_{\cal{C}}$ as a functor in $E$, from the category of
finite sets into the category $\cal{C}$, and this functor is
\emph{exact}, hence transforms finite groups into $\cal{C}$-groups,
etc. If therefore $X$ is an object of~$\cal{C}$ on which the finite group
$G$ operates on the right, we see that one may consider $X$
as an object of~$\cal{C}$ having a $\cal{C}$-group of operators
on the right~$G_{\cal{C}}$. We shall therefore say, by extension of a
general terminology relative to objects with
$\cal{C}$-groups of operators, that $X$ is \emph{formally
principal homogeneous}
under $G$ if $X$ is formally principal homogeneous under
$G_{\cal{C}}$, \ie if the canonical morphism
$$
X\times G_{\cal{C}} \to X\times X
$$
deduced from the operation of~$G_{\cal{C}}$ on~$X$ on the right,
is an isomorphism. We say that $X$ is \emph{principal homogeneous}
under $G$ if it is so under $G_{\cal{C}}$, \ie if it is so formally,
and if moreover the quotient $X/G=X/G_{\cal{C}}$ is~$e_{\cal{C}}$. If one
fixes a fundamental functor, whence an equivalence of
$\cal{C}$ with a category $\cal{C}(\pi)$, $X$ corresponds to
a set on which $\pi$ operates on the left continuously,
say $E=F(X)$. Making $G$ operate on~$X$ on the right then amounts
to making $G$ operate on the set $E$ on the right, in such a way that the operations of~$G$ commute with those of~$\pi$. One
then sees at once that $X$ is principal homogeneous under
$G$ if and only if the set $E$ is a principal
homogeneous space
under $G$,
\ie if and only if $G$ operates on it in a simply transitive way. (Moreover $X$ is formally
principal homogeneous if and only if
$E$ is principal homogeneous \emph{or} empty). Comparing
with~\Ref{V.5.3}, we see that if $X$ is principal homogeneous under
$G$ \emph{and connected}, then the given homomorphism from~$G$ into the
group opposite to $\Aut(X)$ is an \emph{isomorphism}; and
moreover, for an object $X$
of~$\cal{C}$
to be connected and principal homogeneous under the group opposite to
$\Aut(X)$, it is necessary and sufficient, with the notation of \No \Ref{V.4}, that it
be isomorphic to a \emph{Galois} $P_i$. In the typical case
$\cal{C}=\cal{C}(\pi)$, this means that $X$ is isomorphic to a
quotient of~$\pi$ by an invariant subgroup.

Let us still suppose given a fundamental functor~$F$. Then the
datum of an
% original p. 134
$X$ principal homogeneous under a finite group $G$ operating on
the right, and of a point $a\in F(X)$, is equivalent to the
datum of a homomorphism from~$\pi$ into the group~$G$. Indeed,
to such a homomorphism one makes correspond the set $E=G$,
making~$\pi$ operate on it on the left by means of the given homomorphism
$\pi\to G$ and the left translations of~$G$, and
making $G$ operate on it on the right by right translation, the
marked point $a$ of~$E$ being the unit element
of~$G$. Thanks to the foregoing, one thus obtains in an
essentially unique way every triple $(X, G, a)$ having the
properties considered above,
since a
principal homogeneous set with marked point under a group $G$ is identified
with the latter. In this way, one has a direct geometric
interpretation of the functor $G\mto \Hom(\pi, G)$ from the
category of finite groups into~$\Ens$, a functor which is
pro-representable by means of~$\pi$, and whose
consideration would therefore furnish another construction of the group
$\pi$ associated with~$F$.
\end{remark}

\section{Exact functors from one Galois category into another}
\label{V.6}

\begin{proposition}
\label{V.6.1}
Let $\cal{C}$, $\cal{C'}$ be two Galois categories,
$H\colon \cal{C}\to \cal{C'}$ a covariant functor, $F'$ a fundamental
functor on~$\cal{C'}$ and $F=F' \circ H$. The following conditions are
equivalent:
\begin{enumerate}
\item[(i)] $H$ is \emph{exact}, \ie left exact and right
exact.
\item[(ii)] $H$ is left exact, transforms finite sums into
finite sums, and epimorphisms into epimorphisms (or again:
objects $\not\approx \emptyset_{\cal{C}}$ into objects $\not\approx
\emptyset_{\cal{C'}}$).
\item[(iii)] $F$ is a fundamental functor on~$\cal{C}$.
\end{enumerate}
\end{proposition}

The implication (i)$\To$(ii) is a general fact about
categories. Moreover the first form given to (ii)
implies the second, as one sees by noting that if $X$ is an object
of~$\cal{C}$, then $X$ is $\not\approx \emptyset_{\cal{C}}$ if and only if the
morphism $X\to e_{\cal{C}}$ is an epimorphism; one will note that
$F$, being assumed left exact, transforms $e_{\cal{C}}$
into~$e_{\cal{C'}}$. The second form given to (ii)
implies~(iii), since $F$, being left exact and hence
pro-representable, falls under~\Ref{V.5.6}
criterion~(iii). Finally (iii) implies~(i), as follows from the
fact that $F$ is exact, and ``conservative'' (\ie satisfies axiom
(G~6) of~\no \Ref{V.4}).

Let then $\Gamma$ be the fundamental groupoid of~$\cal{C}$, and
$\Gamma'$ that of~$\cal{C'}$. Thus if
% original p. 135
$H$ is exact, then $F' \mto F' \circ H$ is a functor from the
groupoid~$\Gamma'$ into the groupoid~$\Gamma$, which we shall
denote by $\tH$:
$$
\tH (F')(X)=F' (H(X))
$$
which can also be written, with the notation $F(X)=E_X (F)$
introduced in~\no \Ref{V.6}:
$$
E_{H(X)} (F')=E_X (\tH(F')).
$$
This last formula shows, taking into account~\Ref{V.5.8}
or~\Ref{V.4.1}, that the exact functor $H$ is determined (up to
unique isomorphism) when one knows the
corresponding functor~$\tH$. Let us fix an~$F'$, and let $F=\tH (F')$;
then $\tH$ defines a homomorphism from~$\pi_{F'} =\Aut(F')$
into $\pi_F =\Aut(F)$:
$$
\tH\colon \pi_{F'} \to \pi_F \quad (F=\tH (F')=F' \circ H).
$$
Moreover the formula above shows (taking into account~\Ref{V.5.8})
that this homomorphism has the following property: for every
finite set $E$ on which $\pi_F$ operates continuously, the group
$\pi_{F'}$ also operates \emph{continuously} by means of
the preceding homomorphism $\pi_{F'} \to
\pi_F$. \hbox{Applying} this to the quotients of~$\pi_F$ by its open
invariant subgroups, one sees that the preceding condition
also means that the homomorphism under consideration is
continuous. Conversely,
let us be given
an object $F$ of~$\Gamma$, an
object $F'$ of~$\Gamma'$ and a continuous homomorphism
$$
u\colon \pi_{F'} \to \pi_F ,
$$
to it there therefore corresponds a functor from~$\cal{C} (\pi)$
into~$\cal{C}(\pi')$, manifestly exact, hence by virtue
of~\Ref{V.4.1} there corresponds to it a functor $H$ from~$\cal{C}$ into
$\cal{C'}$ which is exact, and such that $\tH\colon \pi_{F'} \to \pi_F$
is precisely~$u$. One can also, instead of a
homomorphism of groups, start from a \emph{functor}
$$
U\colon \Gamma' \to \Gamma
$$
which is such that for \emph{every} $F' \in \Gamma'$ (or for \emph{one} $F'
\in \Gamma'$, which amounts to the same) the corresponding homomorphism
$\pi_{F'} \to \pi_F$ is continuous: such a functor is isomorphic to
a functor of the form~$\tH$, where $H\colon \cal{C} \to
\cal{C'}$ is an exact functor determined up to unique
isomorphism. Thus:

\begin{corollary}
\label{V.6.2}
For
% original p. 136
a functor $H\colon \cal{C} \to \cal{C'}$ between Galois categories
to be exact, it is necessary and sufficient that there exist
equivalences $\cal{C}(\pi)\to\cal{C}$ and
$\cal{C'}\to\cal{C}(\pi')$
which transform the functor $H$ into the functor
$\cal{C}(\pi)\to\cal{C}(\pi')$ associated with a homomorphism of
topological groups $\pi'\to\pi$.
\end{corollary}

\begin{corollary}
\label{V.6.3}
Let $\cal{C}$, $\cal{C'}$ be two Galois categories, and
$\Gamma$, $\Gamma'$ their fundamental groupoids. Then the
category of exact functors from~$\cal{C}$ into $\cal{C'}$ is
equivalent to the category of functors $U\colon \Gamma'
\to \Gamma$ having the following property: for every $F'$ in~$\Gamma'$ (or for \emph{one} $F'$ in~$\Gamma'$, which amounts to the
same), setting $F=U(F')$, the homomorphism
$$
\pi_{F'} =\Aut (F')\to \pi_F =\Aut (F)
$$
defined by $U$ is continuous.
\end{corollary}

Consider the fundamental pro-group $\Pi$ of~$\cal{C}$; then an
exact functor $H$ transforms it into a pro-group $H(\Pi)$
of~$\cal{C'}$; we are going to define a homomorphism
$$
\Pi' \to H(\Pi)
$$
(where $\Pi'$ is the fundamental pro-group of~$\cal{C'}$), by the
condition that for every object $F'$ of~$\Gamma'$, the corresponding
homomorphism
$$
F' (\Pi')=\pi_{F'} \to F' (H(\Pi))=\pi_F \qquad
\text{(where $F=F'\circ H=\tH(F')$)}
$$
be the natural homomorphism
$$
\Aut(F')\to \Aut(F' \circ H).
$$
Since the latter is functorial in~$F'$, it indeed defines, by virtue
of~\Ref{V.5.8}, a homomorphism of pro-objects, and in fact of
pro-groups, of~$\cal{C'}$. This homomorphism is said to be
\emph{associated} with the functor~$H$.

Now let $H'$ be a second exact functor, from the
Galois category $\cal{C'}$ into a
Galois category~$\cal{C''}$. It is trivial that one has
$$
{}^t (H' H)=\tH\,\tH'
$$
(N.B. we have here an identity of functors, and not merely a
canonical isomorphism). One has an analogous transitivity
property for the associated homomorphisms of the
fundamental pro-groups.

We
% original p. 137
shall now interpret the properties of the exact
functor $H$ in terms of the corresponding homomorphism
$$
u\colon \pi_{F'} \to \pi_{F} \qquad \text{(where $F' =F'\circ H$)}.
$$
It is convenient to introduce the notion of a \emph{pointed object}
of the Galois category $\cal{C}$ (endowed with its fundamental
functor $F$): it is by definition an object $X$ of~$\cal{C}$
endowed with an element $a$ of~$F(X)$. It is therefore interpreted
as a finite set on which $\pi_F$ operates continuously on the
left, endowed with a point~$a$. Hence the \emph{connected} pointed
objects of~$\cal{C}$ are identified, by virtue of~\Ref{V.5.3}, with the
open subgroups of~$\pi_F$. If $U$ and $V$ are two such
subgroups, corresponding to connected pointed objects
$X$, $Y$ of~$\cal{C}$, then there exists a pointed homomorphism from
$X$ into $Y$ if and only if $U\subset V$, and this homomorphism is
then unique. Of course, the functor $H$ transforms pointed
objects into pointed objects (since $F=F' \circ H$). Let us note
on the other hand that a closed subgroup of a group such as $\pi_F$
is the intersection of the open subgroups that contain it;
consequently,
if
$M$, $N$ are two closed subgroups, then $M\subset N$ if and
only if every open subgroup that contains $N$ also
contains~$M$. Thanks to these remarks, one easily proves
the results that follow:

\begin{proposition}
\label{V.6.4}
Let $X$ be a connected pointed object of~$\cal{C}$, associated with
an open subgroup $U$ of~$\pi_F$. In order that
$U$
contain $u(\pi_{F'})$ (\resp the closed invariant subgroup
generated by~$u(\pi_{F'})$), it is necessary and sufficient that $H(X)$
admit a pointed section (\resp be completely
decomposed).
\end{proposition}

We shall call a \emph{section} --- understood: over the final
object --- of an object $X$ of a Galois category~$\cal{C}$ a
morphism from the final object $e_{\cal{C}}$ into~$X$, which amounts
to giving an element $a$ of~$F(X)$ invariant under
$\pi_F$; if $X$ is pointed, one says that one has a \emph{pointed
section} if it is compatible with the pointed structures
on~$X$ and~$e_{\cal{C}}$, \ie if $a$ is precisely the marked
object of~$F(X)$. Such a section is therefore unique, and exists if
and only if the marked object of~$F(X)$ is invariant
under~$\pi_F$. Finally, an object of a Galois category is said to be
\emph{completely decomposed}
if it is isomorphic to a sum of final objects, \ie if $\pi_F$
operates trivially on $F(X)$ --- a condition evidently
stronger than the existence of a pointed section, when $X$ is
pointed. Proposition~\Ref{V.6.4} follows trivially from the
preceding definitions and remarks.

\begin{corollary}
\label{V.6.5}
For
% original p. 138
$u$ to be trivial, it is necessary and sufficient that for every object
$X$ of~$\cal{C}$, $H(X)$ be completely decomposed.
\end{corollary}

\begin{proposition}
\label{V.6.6}
Let $X'$ be a connected pointed object of~${\cal{C}}'$, associated
with an open subgroup $U'$ of~$\pi_{F'}$. In order that $U'$
contain~$\Ker u$, it is necessary and sufficient that there exist a
connected pointed object $X$ of~$\cal{C}$ and a pointed homomorphism
from the pointed connected component $X'_0$ of~$H(X)$ into~$X'$ (hence
that $X$ be isomorphic as a pointed object to a quotient of the
neutral connected component of the inverse image of a pointed object
of~$\cal{C}$). If $u$ is surjective, the preceding condition
is also equivalent to the following: $X$ is isomorphic to
an~$H(X)$, where $X$ is a pointed object of~$\cal{C}$.
\end{proposition}

(One calls \emph{neutral connected component}
of a pointed object $X$ of a Galois category~$\cal{C}$
the unique connected pointed subobject of~$X$; it corresponds to the
orbit under $\pi_F$ of the marked point of~$F(X)$, by virtue
of~\Ref{V.5.3}). Since the fact that $U'$ contains $\Ker u$ does not
depend on the chosen pointing of~$X'$ (because another
pointing amounts to replacing $U$ by a subgroup
conjugate to~$U$), one sees:

\begin{corollary}
\label{V.6.7}
In order that $U'$ contain~$\Ker u$, it is necessary and sufficient that there exist an
\hbox{object~$X$} of~$\cal{C}$ (which one may assume connected) and a morphism
from a connected component of~$H(X)$ into~$X'$. If $u$ is surjective,
this also means that $X'$ is isomorphic to an object of the
form~$H(X)$.
\end{corollary}

\begin{corollary}
\label{V.6.8}
In order that $u$ be injective, it is necessary and sufficient that for every object
$X'$ of~${\cal{C}}'$ there exist an object $X$ of~$\cal{C}$ and a
homomorphism from a connected component of~$H(X)$ into~$X'$.
\end{corollary}

\begin{proposition}
\label{V.6.9}
The following conditions are equivalent:
\begin{enumerate}
\item[(i)] The homomorphism $u\colon \pi_{F'} \to \pi_F$ is surjective.
\item[(ii)] For every connected object $X$ of~$\cal{C}$, $H(X)$ is
connected.
\item[(iii)] The functor $H$ is fully faithful.
\end{enumerate}
\end{proposition}

This last fact means that for two objects~$X$, $Y$ of~$\cal{C}$,
the natural map
$$
\Hom(X,Y) \to \Hom(H(X),H(Y))
$$
is bijective.

\begin{corollary}
\label{V.6.10}
For
% original p. 139
$u$ to be an isomorphism, it is necessary and sufficient that $H$ be an
equivalence of categories, or again that the following two conditions
be satisfied:
\begin{enumerate}
\item[a)] for every connected object $X$ of~$\cal{C}$, $H(X)$ is connected
\item[b)] every object of~$\cal{C}'$ is isomorphic to an object of the
form~$H(X)$.
\end{enumerate}
\end{corollary}

\begin{proposition}
\label{V.6.11}
Let $H\colon \cal{C}\to \cal{C}'$ and $H'\colon\cal{C}\to \cal{C}''$
be exact functors between Galois categories, and $F''$ a
fundamental functor on~$\cal{C}''$; let us set $F'=F''H'$ and $F=F'H$, and
consider the associated homomorphisms
$$
u'\colon\pi_{F''}\to \pi_{F'}\qquad u\colon\pi_{F'}\to \pi_F.
$$
In order that $\Ker u\subset \Im u'$, \ie in order that $uu'$ be
the trivial homomorphism, it is necessary and sufficient that for every object $X$
of~$\cal{C}$, $H'(H(X))$ be completely decomposed. In order
that $\Ker u \supset \Im u' $, it is necessary and sufficient that for every
connected pointed object $X'$ of~$\cal{C}'$ such that $H'(X')$ admits a
pointed section, there exist an object~$X$ of~$\cal{C}$ and a
homomorphism from a connected component of~$H(X)$ into~$X'$.
\end{proposition}

The first assertion follows from the last statement
of~\Ref{V.6.4}. The second follows from the conjunction
of~\Ref{V.6.4} and~\Ref{V.6.6}.

\begin{remark}
\label{V.6.12}
It is not true in general, under the conditions
of~\Ref{V.6.8}, that $X'$ is isomorphic to an object of the form
$H(X)$. One can show that in order that every connected object (hence every
object) of~$\cal{C}'$ be isomorphic to an object of the form $H(X)$,
it is necessary and sufficient that $u$ be an isomorphism of~$\pi_{F'}$ onto a
subgroup of~$\pi_F$ that is a \emph{direct factor}. In practice, however,
one directly constructs a homomorphism $\pi_F\to\pi_{F'}$ right inverse
to~$u$ by means of a suitable exact functor
from~$\cal{C'}$ into~$\cal{C}$.
\end{remark}

\begin{proposition}
\label{V.6.13}
Let $\cal{C}$ be a Galois category endowed with a fundamental
functor $F$, $S$ a connected object of~$\cal{C}$, and $\cal{C}'$ the
category of objects of~$\cal{C}$ over~$S$. Then
$\cal{C'}$ is a Galois category, and the functor $X\mto
H(X)=X\times S$ from~$\cal{C}$ into~$\cal{C}'$ is exact. Let $a\in
F(S)$, and let $F'$ be the functor from~$\cal{C}'$ into the category
of finite sets defined by
$$
F'(X')=\text{inverse image of~$a$ under }F(X')\to F(S).
$$
Then one has an isomorphism $F\cong F'\circ H$, and the
corresponding homomorphism
$$
u:\pi_{F'}\to \pi_F
$$
% original p. 140
is an isomorphism of~$\pi_{F'}$ onto the open subgroup $U$ of
$\pi_{F}$ that is the stabilizer of the marked element $a$ of~$F(X)$.
\end{proposition}

The proof is left to the reader.

\section{The case of preschemes}
\label{V.7}
Let $S$ be a locally noetherian and
\emph{connected} prescheme, and let
$$
a\colon \Spec(\Omega)\to S
$$
be a geometric point of~$S$, with values in an
algebraically closed field $\Omega$. We shall set
$$
\cal{C}=\text{category of \'etale coverings of~$S$}
$$
and for an object $X$ of~$\cal{C}$, \ie an \'etale covering
$X$ of~$S$, we set
$$
F(X)=\text{set of geometric points of~$X$
over~$a$.}
$$
Thus $F$ becomes a functor on~$\cal{C}$ with values in the
category of finite sets. Properties (G~1) to
(G~6) are satisfied: for (G~1), this is contained in the sorites of
I~\Ref{I.4.6}, (G~2) follows from~\Ref{V.3.4}, (G~3) from~\Ref{V.3.5},
(G~4) is trivial by definition, (G~5) follows from~\Ref{V.3.5} and
from the beginning of \No \Ref{V.2}, finally (G~6) is proved in~\Ref{V.3.7}. One
can therefore apply the results of \No \Ref{V.4}, \Ref{V.5}, \Ref{V.6}. This makes it
possible in particular to define a pro-object $P$ of~$\cal{C}$
representing $F$, called the \emph{universal covering of~$S$
at the point~$a$},
and a topological group $\pi=\Aut(F)=\Aut^0 P$, called the
\emph{fundamental group of~$S$ at~$a$},
and denoted~$\pi_1(S,a)$.
The functor $F$ then defines an equivalence of the
category $\cal{C}$ with the category of finite sets
on which $\pi=\pi_1(S,a)$ operates continuously. This
equivalence therefore makes it possible to interpret the usual
operations of projective limits and finite inductive limits on
coverings (products, fiber products, sums, passage to the
quotient, etc.) in terms of the analogous operations in
$\cal{C}(\pi)$, \ie in terms of the obvious operations on
finite sets on which $\pi$ operates. Let us note moreover, since
the topological connected components of an \'etale covering
are also \'etale coverings, that \emph{an \hbox{object $X$} of~$\cal{C}$ is connected in $\cal{C}$ if and only if it is
topologically connected}; by virtue of~\Ref{V.5.3} this therefore means
that $\pi_1$ operates transitively on~$F(X)$.
% original p. 141
Let us note that in order for an object $X$ of~$\cal{C}$ to be faithfully flat
and quasi-compact over~$S$ (as it is already flat and
quasi-compact over~$S$), it is necessary and sufficient that $X\to S$ be surjective,
\ie be an epimorphism in $\cal{C}$, or again that $X\neq
\emptyset$. One then concludes from criterion \Ref{V.2.6}~(iii) that~$X$
\emph{is a principal covering of
$S$
with group $G$ if and only if it is a principal homogeneous space
under $G$ in the category $\cal{C}$}, (as it was
defined in \No \Ref{V.5}).

If $a'$ is another geometric point of~$S$ (corresponding
to an algebraically closed field $\Omega'$, which may be
different from~$\Omega$ and which may even have a
different characteristic), it defines a fiber functor
$F'=F_{a'}$ from~$\cal{C}$ into the category of finite sets,
which is again exact, hence isomorphic to $F=F_a$. Consequently, the
fundamental groups $\pi_1(S;a)$
for variable $a$ are isomorphic to one another. If one denotes by
$\pi_1(S;a,a')$
the set of isomorphisms (or, what amounts to the same,
the set of homomorphisms) $F_a\to F_{a'}$ of the associated fiber functors,
one thus obtains a \emph{groupoid} whose
set of objects is the set of geometric points
of~$S$, the fundamental groups being the automorphism
groups of the objects of the said groupoid. The set
$\pi_1(S;a',a)$ may be called the \emph{set of classes
of paths from~$a$ to $a'$}. These classes therefore compose in the obvious way. Finally, one can define a pro-group $\Pi_1^S$
of~$\cal{C}$, which one may call the \emph{fundamental pro-group
of}~$S$
or \emph{local system of fundamental groups on}~$S$,
defined up to unique isomorphism by the condition that one
have an isomorphism, functorial in the geometric point $a$ of
$S$:
$$
F_a(\Pi_1^S)=\pi_1(S;a)
$$
(\cf remark~\Ref{V.5.10}). In particular, if $s$ is an ordinary
point of~$S$, the fiber of~$G$ at~$s$ is a pro-group over~$\kres(s)$,
a projective limit of finite \'etale groups over~$\kres(s)$; one could
call this pro-group \emph{the fundamental group of~$S$ at the
ordinary point $s$ of~$S$}, and denote it $\pi_1(S,s)$. By definition, these
points with values in an algebraically closed extension
$\Omega$ of~$\kres(s)$ are the elements of~$\pi_1(S;a)$, where
$a$ is the geometric point of~$S$ defined by the said
extension. In particular, (taking for $S$ the spectrum of a field)
to every field $k$ there is associated canonically and functorially
a pro-group over~$k$, which one could denote $\pi_1(k)$, a
projective limit of finite \'etale groups over~$k$, and whose points
in an algebraically closed extension $\Omega$ of~$k$
are identified
% original p. 142
with the elements of the topological Galois group of~$\bar{k}/k$,
where $\bar{k}$ is the Galois closure of~$k$ in $\Omega$
(\cf \Ref{V.8.1}). This group $\pi_1(k)$ does not yet seem to have
attracted the attention of algebraists.

Now let
$$
f\colon S'\to S
$$
be a morphism from one connected locally noetherian prescheme
into another, let $a'$ be a geometric point of~$S'$ and let
$a=f(a')$ be its direct image in $S$. Then the ``inverse
image'' functor induces a functor from the category $\cal{C}(S)$
of \'etale coverings of~$S$ into the category
$\cal{C}(S')$ of \'etale coverings of~$S'$:
$$
f^\bbullet:\cal{C}(S)\to \cal{C}(S')
$$
One has moreover an isomorphism of functors
$$
F_a\cong F_{a'}\circ f^\bbullet,
$$
so that $f^\bbullet$ is an \emph{exact} functor, to which
the results of \No \Ref{V.6} apply. In particular one has a
canonical homomorphism
$$
u=\pi_1(f;a'):\pi_1(S',a')\to \pi_1(S;a) \qquad (a'=f(a))
$$
which makes it possible to reconstruct the inverse image functor, as an
operation of restriction of the groups of operators. The
properties of the functor $f^\bbullet$ are therefore expressed in a
simple way by the properties of the associated homomorphism of
groups, as was made explicit in
\No \Ref{V.6}. If in particular $S'$ is an \'etale covering of~$S$,
then the homomorphism $u$ is an isomorphism of~$\pi_1(S',a')$ onto the
open subgroup of~$\pi_1(S,a)$ that defines the pointed connected
\'etale covering $S'$ of~$S$ (\ie the stabilizer~$U$
of~$a'\in F_a(S')$ in $\pi_1(S;A)$).

If one wishes to interpret the homomorphisms $\pi_1(f;a')$ for
a variable geometric point $a'$, one must, in accordance
with what was said in \No \Ref{V.6}, consider a
homomorphism
$$
\Pi_1(f):\Pi_1^{S'}\to f^\bbullet(\Pi_1^S)
$$
of pro-groups over~$S$, and take the corresponding homomorphism for
the geometric fibers.

\section{The case of a normal base prescheme}
\label{V.8}
% original p. 143
\begin{proposition}
\label{V.8.1}
Let $S$ be the spectrum of a field $k$, and let $\Omega$ be an
algebraically closed extension of~$k$, defining a geometric
point $a$ of~$S$ with values in $\Omega$. Let
$\bar{k}$ be the separable closure of~$k$ in $\Omega$. Then there
exists a canonical isomorphism of~$\pi_1(S,a)$ onto the topological
Galois group of~$\bar{k}/k$.
\end{proposition}

Let $k'$ be the algebraic closure of~$k$ in $\Omega$; there
therefore corresponds to it a geometric point~$b$ of~$S$, with
values in $k'$. The natural homomorphism of functors $F_b\to F_a$
is evidently an isomorphism, since a $k$-homomorphism of a
finite separable extension of~$k$ into $\Omega$
necessarily takes its values in $\bar{k}$ and a fortiori in
$k'$. On the other hand, the group $\pi'$ of $k$-automorphisms of~$k'/k$
operates in the obvious way on~$F_b$, whence a
homomorphism $\pi'\to \Aut(F_b)\isomto \Aut(F_a)=\pi_1(S;a)$.
On the other hand, it is well known that the natural homomorphism from~$\pi'$
into the group $\pi$ of automorphisms of~$\bar{k}/k$ is an
isomorphism. One thus obtains a canonical homomorphism $\pi\to
\pi_1(S;a)$; it remains to show that it is an isomorphism. Indeed,
this homomorphism is injective, since an element of the kernel is an
automorphism of~$\bar{k}/k$ which induces the identity on every
finite separable subextension, hence is trivial. This
homomorphism is surjective, since if $X$ is a \emph{connected}
\'etale covering of~$S$, hence defined by a finite separable
\emph{extension}~$L/k$, then $\pi$ is transitive
on the set of $k$-homomorphisms of~$L$ into $k'$, as is well
known.

\begin{proposition}
\label{V.8.2}
Let $S$ be a connected, locally noetherian and
normal prescheme, $K=\kres(s)$ its function field $=$ the residue field at
its generic point $s$, and $\Omega$ an
algebraically closed extension of~$K$, defining a geometric
point $a'$ of~$S'=\Spec(K)$ and a geometric
point $a$ of~$S$. Then the homomorphism
$\pi_1(S';a')\to\pi_1(S;a)$ is surjective. When one identifies the
first group with the Galois group of the separable closure
$\bar{K}$ of~$K$ in~$\Omega$ (\cf \Ref{V.8.1}), then the kernel of
the preceding homomorphism corresponds by Galois theory
to the subextension of~$\bar{K}/K$ composed of the
finite extensions of~$K$ in $\Omega$ which are unramified over~$S$.
\end{proposition}

The first assertion means that the inverse image on~$S'$ of a
connected \'etale covering $X$ of~$S$ is connected, \ie that $X$
is integral; this is none other than (I~\Ref{I.10.1}). The kernel of
the preceding homomorphism is then interpreted as
consisting of the automorphisms of~$\bar{K}/K$ which induce
the identity on the sets~$F_a(X)$,
% original p. 144
where one may assume the \'etale covering $X$ of~$S$
connected. But this means that this automorphism induces
the identity on the finite subextensions of~$\bar{K}/K$ which are
unramified over~$S$, which proves the last assertion.

\begin{remarkstar}
Thanks to this interpretation of the fundamental group of the
normal prescheme $S$ in terms of the usual Galois theory,
the definition had been known in this case for a
long time.
\end{remarkstar}

\section{The case of non-connected preschemes: multi-Galois
categories}
\label{V.9}

Let $S$ be a locally noetherian prescheme, and let
$(S_i)_{i\in I}$ be its connected components. Then the category
$\cal{C}(S)$ of \'etale coverings of~$S$ is equivalent
to the product category of the $\cal{C}(S_i)$, which
are interpreted in terms of the fundamental groups of the $S_i$, once
a geometric point has been chosen in each $S_i$. In the
application of descent theory for \'etale morphisms, it is
sometimes inconvenient to choose, for every $S_i$,
a geometric point of~$S_i$. It is then more convenient
to resort to the natural generalization of~\Ref{V.5.8}
in order to interpret $\cal{C}(S)$ as a category of functors
on the groupoid of geometric points of~$S$,
considered as the sum of the groupoids corresponding to the
connected components of~$S$; the functors in question are the
functors with values in the category of finite sets
satisfying the continuity property analogous to
the one invoked in~\Ref{V.5.8}. In practice, one will have a family
$(a_t)_{t\in E}$ of geometric points of~$S$ such that
every connected component $S_i$ of~$S$ contains at least one of them, and one
can then, as in~\Ref{V.5.10}, replace the groupoid of
all geometric points of~$S$ by the
analogous groupoid whose underlying set is~$E$. Of course, these
considerations ought to fit into
general definitions concerning the categories which are
equivalent to product categories of categories
of the form $\cal{C}(\pi)$, and which one may call
\emph{multi-Galois categories}.
We shall leave the details to the reader.

\end{document}
